\documentclass[10pt,onecolumn]{IEEEtran}
\usepackage{amsmath,amssymb,booktabs,longtable,array,pdflscape,url,hyperref}
\hypersetup{hidelinks}
\newcommand{\Hnone}{\textsc{H-None}}
\newcommand{\Hsummary}{\textsc{H-Summary}}
\newcommand{\Hfull}{\textsc{H-Full}}
\begin{document}
\title{Supplementary Material: NetShield}
\author{Nuonan Ouyang, Adrian Shatte, Zhigang Lu, Chao Chen, and Wei Xiang}
\maketitle

\section{Evidence Classification and Scope}
The formal Raspberry Pi campaign contains 12 sessions and 216 accepted arms. Task C reconstructs timing from accepted formal hardware timestamps. Task A is offline replay, Task B is a software sensitivity, and Task D is discrete-event conformance. None is a new Raspberry Pi experiment. This supplementary build does not modify the frozen formal raw evidence.

\begin{table}[h]\centering\caption{Evidence classes.}\begin{tabular}{lll}\toprule
Component & Evidence type & Permitted interpretation\\\midrule
Formal campaign & Hardware & Three-Pi measurements\\
Full-raw audit & Raw reconstruction & Accepted-arm traceability\\
Task A & Offline replay & Frozen-selector sensitivity\\
Task B & Software sensitivity & Phase-desynchronization counterfactual\\
Task C & Raw reconstruction & Accepted-hardware timestamp metrics\\
Task D & Discrete-event conformance & Analytical boundary check\\\bottomrule
\end{tabular}\end{table}

\section{Frozen Action Table and I-QPG Thresholds}
\begin{table}[h]\centering\caption{Frozen action lattice and $v=1$ thresholds against H-NONE.}\begin{tabular}{lrrr}\toprule
Action & Payload (B) & Utility & Threshold $Q$ (B)\\\midrule
L-NONE&0&0.758501&--\\ L-ALERT&256&0.825168&221.84\\ L-SUMMARY&4096&0.891835&30.14\\
M-NONE&0&0.766968&--\\ M-ALERT&256&0.833635&254.91\\ M-SUMMARY&8192&0.900302&16.10\\
H-NONE&0&0.768378&--\\ H-ALERT&256&0.835044&260.41\\ H-SUMMARY&16384&0.901711&8.14\\ H-FULL&131072&0.968378&1.53\\\bottomrule
\end{tabular}\end{table}

\section{Theorem Derivations}
For $m$ remotely represented probes and $k$ full-evidence probes, payload is $mS+k(B-S)$. Hence one-round feasibility is
\[kB+(m-k)S\le C,\]
and, for $C\ge mS$,
\[k^*(C,m)=\min\left\{m,\left\lfloor\frac{C-mS}{B-S}\right\rfloor\right\}.\]
With $\rho=C/(nB)$, $\sigma=S/B$, $r=m/n$, and $f=k/n$, the continuous frontier is $\sigma r+(1-\sigma)f\le\rho$. For synchronized fixed-full overload, conservation gives $Q_\Sigma(T)=T(nB-C)$. At mid capacity the frozen I-QPG state alternates between H-FULL and H-NONE because the post-full queue exceeds every positive-payload threshold and the zero-payload round drains the residual.

\begin{landscape}
\section{Complete V-Sweep Results}
The $v=1$ identity gate covers 36 accepted D cells and four deterministic checks per cell. Low-capacity cycle/action composition changes with $v$; no all-$v$ cycle-identity claim is made.
\scriptsize\begin{longtable}{lllr r r r}
\caption{Complete frozen-selector $v$-sweep results (432 retained rows).}\label{tab:vsweepfull}\\
\toprule
Session & $v$ & Cap. & Core frac. & Payload (B) & Max Q (B) & Period \\
\midrule
\endfirsthead
\toprule
Session & $v$ & Cap. & Core frac. & Payload (B) & Max Q (B) & Period \\
\midrule
\endhead
F12\_S01 & 1 & low & 0.244444 & 9191424 & 87382 & 4 \\
F12\_S02 & 1 & low & 0.244444 & 9191424 & 87382 & 4 \\
F12\_S03 & 1 & low & 0.244444 & 9191424 & 87382 & 4 \\
F12\_S04 & 1 & low & 0.244444 & 9191424 & 87382 & 4 \\
F12\_S05 & 1 & low & 0.244444 & 9191424 & 87382 & 4 \\
F12\_S06 & 1 & low & 0.244444 & 9191424 & 87382 & 4 \\
F12\_S07 & 1 & low & 0.244444 & 9191424 & 87382 & 4 \\
F12\_S08 & 1 & low & 0.244444 & 9191424 & 87382 & 4 \\
F12\_S09 & 1 & low & 0.244444 & 9191424 & 87382 & 4 \\
F12\_S10 & 1 & low & 0.244444 & 9191424 & 87382 & 4 \\
F12\_S11 & 1 & low & 0.244444 & 9191424 & 87382 & 4 \\
F12\_S12 & 1 & low & 0.244444 & 9191424 & 87382 & 4 \\
F12\_S01 & 1 & mid & 0.500000 & 17694720 & 43691 & 2 \\
F12\_S02 & 1 & mid & 0.500000 & 17694720 & 43691 & 2 \\
F12\_S03 & 1 & mid & 0.500000 & 17694720 & 43691 & 2 \\
F12\_S04 & 1 & mid & 0.500000 & 17694720 & 43691 & 2 \\
F12\_S05 & 1 & mid & 0.500000 & 17694720 & 43691 & 2 \\
F12\_S06 & 1 & mid & 0.500000 & 17694720 & 43691 & 2 \\
F12\_S07 & 1 & mid & 0.500000 & 17694720 & 43691 & 2 \\
F12\_S08 & 1 & mid & 0.500000 & 17694720 & 43691 & 2 \\
F12\_S09 & 1 & mid & 0.500000 & 17694720 & 43691 & 2 \\
F12\_S10 & 1 & mid & 0.500000 & 17694720 & 43691 & 2 \\
F12\_S11 & 1 & mid & 0.500000 & 17694720 & 43691 & 2 \\
F12\_S12 & 1 & mid & 0.500000 & 17694720 & 43691 & 2 \\
F12\_S01 & 1 & high & 1.000000 & 35389440 & 0 & 1 \\
F12\_S02 & 1 & high & 1.000000 & 35389440 & 0 & 1 \\
F12\_S03 & 1 & high & 1.000000 & 35389440 & 0 & 1 \\
F12\_S04 & 1 & high & 1.000000 & 35389440 & 0 & 1 \\
F12\_S05 & 1 & high & 1.000000 & 35389440 & 0 & 1 \\
F12\_S06 & 1 & high & 1.000000 & 35389440 & 0 & 1 \\
F12\_S07 & 1 & high & 1.000000 & 35389440 & 0 & 1 \\
F12\_S08 & 1 & high & 1.000000 & 35389440 & 0 & 1 \\
F12\_S09 & 1 & high & 1.000000 & 35389440 & 0 & 1 \\
F12\_S10 & 1 & high & 1.000000 & 35389440 & 0 & 1 \\
F12\_S11 & 1 & high & 1.000000 & 35389440 & 0 & 1 \\
F12\_S12 & 1 & high & 1.000000 & 35389440 & 0 & 1 \\
F12\_S01 & 3 & low & 0.244444 & 9191424 & 87382 & 4 \\
F12\_S02 & 3 & low & 0.244444 & 9191424 & 87382 & 4 \\
F12\_S03 & 3 & low & 0.244444 & 9191424 & 87382 & 4 \\
F12\_S04 & 3 & low & 0.244444 & 9191424 & 87382 & 4 \\
F12\_S05 & 3 & low & 0.244444 & 9191424 & 87382 & 4 \\
F12\_S06 & 3 & low & 0.244444 & 9191424 & 87382 & 4 \\
F12\_S07 & 3 & low & 0.244444 & 9191424 & 87382 & 4 \\
F12\_S08 & 3 & low & 0.244444 & 9191424 & 87382 & 4 \\
F12\_S09 & 3 & low & 0.244444 & 9191424 & 87382 & 4 \\
F12\_S10 & 3 & low & 0.244444 & 9191424 & 87382 & 4 \\
F12\_S11 & 3 & low & 0.244444 & 9191424 & 87382 & 4 \\
F12\_S12 & 3 & low & 0.244444 & 9191424 & 87382 & 4 \\
F12\_S01 & 3 & mid & 0.500000 & 17694720 & 43691 & 2 \\
F12\_S02 & 3 & mid & 0.500000 & 17694720 & 43691 & 2 \\
F12\_S03 & 3 & mid & 0.500000 & 17694720 & 43691 & 2 \\
F12\_S04 & 3 & mid & 0.500000 & 17694720 & 43691 & 2 \\
F12\_S05 & 3 & mid & 0.500000 & 17694720 & 43691 & 2 \\
F12\_S06 & 3 & mid & 0.500000 & 17694720 & 43691 & 2 \\
F12\_S07 & 3 & mid & 0.500000 & 17694720 & 43691 & 2 \\
F12\_S08 & 3 & mid & 0.500000 & 17694720 & 43691 & 2 \\
F12\_S09 & 3 & mid & 0.500000 & 17694720 & 43691 & 2 \\
F12\_S10 & 3 & mid & 0.500000 & 17694720 & 43691 & 2 \\
F12\_S11 & 3 & mid & 0.500000 & 17694720 & 43691 & 2 \\
F12\_S12 & 3 & mid & 0.500000 & 17694720 & 43691 & 2 \\
F12\_S01 & 3 & high & 1.000000 & 35389440 & 0 & 1 \\
F12\_S02 & 3 & high & 1.000000 & 35389440 & 0 & 1 \\
F12\_S03 & 3 & high & 1.000000 & 35389440 & 0 & 1 \\
F12\_S04 & 3 & high & 1.000000 & 35389440 & 0 & 1 \\
F12\_S05 & 3 & high & 1.000000 & 35389440 & 0 & 1 \\
F12\_S06 & 3 & high & 1.000000 & 35389440 & 0 & 1 \\
F12\_S07 & 3 & high & 1.000000 & 35389440 & 0 & 1 \\
F12\_S08 & 3 & high & 1.000000 & 35389440 & 0 & 1 \\
F12\_S09 & 3 & high & 1.000000 & 35389440 & 0 & 1 \\
F12\_S10 & 3 & high & 1.000000 & 35389440 & 0 & 1 \\
F12\_S11 & 3 & high & 1.000000 & 35389440 & 0 & 1 \\
F12\_S12 & 3 & high & 1.000000 & 35389440 & 0 & 1 \\
F12\_S01 & 10 & low & 0.300000 & 10838016 & 87386 & 10 \\
F12\_S02 & 10 & low & 0.300000 & 10838016 & 87386 & 10 \\
F12\_S03 & 10 & low & 0.300000 & 10838016 & 87386 & 10 \\
F12\_S04 & 10 & low & 0.300000 & 10838016 & 87386 & 10 \\
F12\_S05 & 10 & low & 0.300000 & 10838016 & 87386 & 10 \\
F12\_S06 & 10 & low & 0.300000 & 10838016 & 87386 & 10 \\
F12\_S07 & 10 & low & 0.300000 & 10838016 & 87386 & 10 \\
F12\_S08 & 10 & low & 0.300000 & 10838016 & 87386 & 10 \\
F12\_S09 & 10 & low & 0.300000 & 10838016 & 87386 & 10 \\
F12\_S10 & 10 & low & 0.300000 & 10838016 & 87386 & 10 \\
F12\_S11 & 10 & low & 0.300000 & 10838016 & 87386 & 10 \\
F12\_S12 & 10 & low & 0.300000 & 10838016 & 87386 & 10 \\
F12\_S01 & 10 & mid & 0.500000 & 17694720 & 43691 & 2 \\
F12\_S02 & 10 & mid & 0.500000 & 17694720 & 43691 & 2 \\
F12\_S03 & 10 & mid & 0.500000 & 17694720 & 43691 & 2 \\
F12\_S04 & 10 & mid & 0.500000 & 17694720 & 43691 & 2 \\
F12\_S05 & 10 & mid & 0.500000 & 17694720 & 43691 & 2 \\
F12\_S06 & 10 & mid & 0.500000 & 17694720 & 43691 & 2 \\
F12\_S07 & 10 & mid & 0.500000 & 17694720 & 43691 & 2 \\
F12\_S08 & 10 & mid & 0.500000 & 17694720 & 43691 & 2 \\
F12\_S09 & 10 & mid & 0.500000 & 17694720 & 43691 & 2 \\
F12\_S10 & 10 & mid & 0.500000 & 17694720 & 43691 & 2 \\
F12\_S11 & 10 & mid & 0.500000 & 17694720 & 43691 & 2 \\
F12\_S12 & 10 & mid & 0.500000 & 17694720 & 43691 & 2 \\
F12\_S01 & 10 & high & 1.000000 & 35389440 & 0 & 1 \\
F12\_S02 & 10 & high & 1.000000 & 35389440 & 0 & 1 \\
F12\_S03 & 10 & high & 1.000000 & 35389440 & 0 & 1 \\
F12\_S04 & 10 & high & 1.000000 & 35389440 & 0 & 1 \\
F12\_S05 & 10 & high & 1.000000 & 35389440 & 0 & 1 \\
F12\_S06 & 10 & high & 1.000000 & 35389440 & 0 & 1 \\
F12\_S07 & 10 & high & 1.000000 & 35389440 & 0 & 1 \\
F12\_S08 & 10 & high & 1.000000 & 35389440 & 0 & 1 \\
F12\_S09 & 10 & high & 1.000000 & 35389440 & 0 & 1 \\
F12\_S10 & 10 & high & 1.000000 & 35389440 & 0 & 1 \\
F12\_S11 & 10 & high & 1.000000 & 35389440 & 0 & 1 \\
F12\_S12 & 10 & high & 1.000000 & 35389440 & 0 & 1 \\
F12\_S01 & 30 & low & 0.322222 & 11476992 & 87398 & 28 \\
F12\_S02 & 30 & low & 0.322222 & 11476992 & 87398 & 28 \\
F12\_S03 & 30 & low & 0.322222 & 11476992 & 87398 & 28 \\
F12\_S04 & 30 & low & 0.322222 & 11476992 & 87398 & 28 \\
F12\_S05 & 30 & low & 0.322222 & 11476992 & 87398 & 28 \\
F12\_S06 & 30 & low & 0.322222 & 11476992 & 87398 & 28 \\
F12\_S07 & 30 & low & 0.322222 & 11476992 & 87398 & 28 \\
F12\_S08 & 30 & low & 0.322222 & 11476992 & 87398 & 28 \\
F12\_S09 & 30 & low & 0.322222 & 11476992 & 87398 & 28 \\
F12\_S10 & 30 & low & 0.322222 & 11476992 & 87398 & 28 \\
F12\_S11 & 30 & low & 0.322222 & 11476992 & 87398 & 28 \\
F12\_S12 & 30 & low & 0.322222 & 11476992 & 87398 & 28 \\
F12\_S01 & 30 & mid & 0.500000 & 17694720 & 43691 & 2 \\
F12\_S02 & 30 & mid & 0.500000 & 17694720 & 43691 & 2 \\
F12\_S03 & 30 & mid & 0.500000 & 17694720 & 43691 & 2 \\
F12\_S04 & 30 & mid & 0.500000 & 17694720 & 43691 & 2 \\
F12\_S05 & 30 & mid & 0.500000 & 17694720 & 43691 & 2 \\
F12\_S06 & 30 & mid & 0.500000 & 17694720 & 43691 & 2 \\
F12\_S07 & 30 & mid & 0.500000 & 17694720 & 43691 & 2 \\
F12\_S08 & 30 & mid & 0.500000 & 17694720 & 43691 & 2 \\
F12\_S09 & 30 & mid & 0.500000 & 17694720 & 43691 & 2 \\
F12\_S10 & 30 & mid & 0.500000 & 17694720 & 43691 & 2 \\
F12\_S11 & 30 & mid & 0.500000 & 17694720 & 43691 & 2 \\
F12\_S12 & 30 & mid & 0.500000 & 17694720 & 43691 & 2 \\
F12\_S01 & 30 & high & 1.000000 & 35389440 & 0 & 1 \\
F12\_S02 & 30 & high & 1.000000 & 35389440 & 0 & 1 \\
F12\_S03 & 30 & high & 1.000000 & 35389440 & 0 & 1 \\
F12\_S04 & 30 & high & 1.000000 & 35389440 & 0 & 1 \\
F12\_S05 & 30 & high & 1.000000 & 35389440 & 0 & 1 \\
F12\_S06 & 30 & high & 1.000000 & 35389440 & 0 & 1 \\
F12\_S07 & 30 & high & 1.000000 & 35389440 & 0 & 1 \\
F12\_S08 & 30 & high & 1.000000 & 35389440 & 0 & 1 \\
F12\_S09 & 30 & high & 1.000000 & 35389440 & 0 & 1 \\
F12\_S10 & 30 & high & 1.000000 & 35389440 & 0 & 1 \\
F12\_S11 & 30 & high & 1.000000 & 35389440 & 0 & 1 \\
F12\_S12 & 30 & high & 1.000000 & 35389440 & 0 & 1 \\
F12\_S01 & 100 & low & 0.333333 & 11821056 & 87436 & -- \\
F12\_S02 & 100 & low & 0.333333 & 11821056 & 87436 & -- \\
F12\_S03 & 100 & low & 0.333333 & 11821056 & 87436 & -- \\
F12\_S04 & 100 & low & 0.333333 & 11821056 & 87436 & -- \\
F12\_S05 & 100 & low & 0.333333 & 11821056 & 87436 & -- \\
F12\_S06 & 100 & low & 0.333333 & 11821056 & 87436 & -- \\
F12\_S07 & 100 & low & 0.333333 & 11821056 & 87436 & -- \\
F12\_S08 & 100 & low & 0.333333 & 11821056 & 87436 & -- \\
F12\_S09 & 100 & low & 0.333333 & 11821056 & 87436 & -- \\
F12\_S10 & 100 & low & 0.333333 & 11821056 & 87436 & -- \\
F12\_S11 & 100 & low & 0.333333 & 11821056 & 87436 & -- \\
F12\_S12 & 100 & low & 0.333333 & 11821056 & 87436 & -- \\
F12\_S01 & 100 & mid & 0.500000 & 17694720 & 43691 & 2 \\
F12\_S02 & 100 & mid & 0.500000 & 17694720 & 43691 & 2 \\
F12\_S03 & 100 & mid & 0.500000 & 17694720 & 43691 & 2 \\
F12\_S04 & 100 & mid & 0.500000 & 17694720 & 43691 & 2 \\
F12\_S05 & 100 & mid & 0.500000 & 17694720 & 43691 & 2 \\
F12\_S06 & 100 & mid & 0.500000 & 17694720 & 43691 & 2 \\
F12\_S07 & 100 & mid & 0.500000 & 17694720 & 43691 & 2 \\
F12\_S08 & 100 & mid & 0.500000 & 17694720 & 43691 & 2 \\
F12\_S09 & 100 & mid & 0.500000 & 17694720 & 43691 & 2 \\
F12\_S10 & 100 & mid & 0.500000 & 17694720 & 43691 & 2 \\
F12\_S11 & 100 & mid & 0.500000 & 17694720 & 43691 & 2 \\
F12\_S12 & 100 & mid & 0.500000 & 17694720 & 43691 & 2 \\
F12\_S01 & 100 & high & 1.000000 & 35389440 & 0 & 1 \\
F12\_S02 & 100 & high & 1.000000 & 35389440 & 0 & 1 \\
F12\_S03 & 100 & high & 1.000000 & 35389440 & 0 & 1 \\
F12\_S04 & 100 & high & 1.000000 & 35389440 & 0 & 1 \\
F12\_S05 & 100 & high & 1.000000 & 35389440 & 0 & 1 \\
F12\_S06 & 100 & high & 1.000000 & 35389440 & 0 & 1 \\
F12\_S07 & 100 & high & 1.000000 & 35389440 & 0 & 1 \\
F12\_S08 & 100 & high & 1.000000 & 35389440 & 0 & 1 \\
F12\_S09 & 100 & high & 1.000000 & 35389440 & 0 & 1 \\
F12\_S10 & 100 & high & 1.000000 & 35389440 & 0 & 1 \\
F12\_S11 & 100 & high & 1.000000 & 35389440 & 0 & 1 \\
F12\_S12 & 100 & high & 1.000000 & 35389440 & 0 & 1 \\
F12\_S01 & 168 & low & 0.244444 & 9209088 & 87382 & 4 \\
F12\_S02 & 168 & low & 0.244444 & 9209088 & 87382 & 4 \\
F12\_S03 & 168 & low & 0.244444 & 9209088 & 87382 & 4 \\
F12\_S04 & 168 & low & 0.244444 & 9209088 & 87382 & 4 \\
F12\_S05 & 168 & low & 0.244444 & 9209088 & 87382 & 4 \\
F12\_S06 & 168 & low & 0.244444 & 9209088 & 87382 & 4 \\
F12\_S07 & 168 & low & 0.244444 & 9209088 & 87382 & 4 \\
F12\_S08 & 168 & low & 0.244444 & 9209088 & 87382 & 4 \\
F12\_S09 & 168 & low & 0.244444 & 9209088 & 87382 & 4 \\
F12\_S10 & 168 & low & 0.244444 & 9209088 & 87382 & 4 \\
F12\_S11 & 168 & low & 0.244444 & 9209088 & 87382 & 4 \\
F12\_S12 & 168 & low & 0.244444 & 9209088 & 87382 & 4 \\
F12\_S01 & 168 & mid & 0.500000 & 17729280 & 43691 & 2 \\
F12\_S02 & 168 & mid & 0.500000 & 17729280 & 43691 & 2 \\
F12\_S03 & 168 & mid & 0.500000 & 17729280 & 43691 & 2 \\
F12\_S04 & 168 & mid & 0.500000 & 17729280 & 43691 & 2 \\
F12\_S05 & 168 & mid & 0.500000 & 17729280 & 43691 & 2 \\
F12\_S06 & 168 & mid & 0.500000 & 17729280 & 43691 & 2 \\
F12\_S07 & 168 & mid & 0.500000 & 17729280 & 43691 & 2 \\
F12\_S08 & 168 & mid & 0.500000 & 17729280 & 43691 & 2 \\
F12\_S09 & 168 & mid & 0.500000 & 17729280 & 43691 & 2 \\
F12\_S10 & 168 & mid & 0.500000 & 17729280 & 43691 & 2 \\
F12\_S11 & 168 & mid & 0.500000 & 17729280 & 43691 & 2 \\
F12\_S12 & 168 & mid & 0.500000 & 17729280 & 43691 & 2 \\
F12\_S01 & 168 & high & 1.000000 & 35389440 & 0 & 1 \\
F12\_S02 & 168 & high & 1.000000 & 35389440 & 0 & 1 \\
F12\_S03 & 168 & high & 1.000000 & 35389440 & 0 & 1 \\
F12\_S04 & 168 & high & 1.000000 & 35389440 & 0 & 1 \\
F12\_S05 & 168 & high & 1.000000 & 35389440 & 0 & 1 \\
F12\_S06 & 168 & high & 1.000000 & 35389440 & 0 & 1 \\
F12\_S07 & 168 & high & 1.000000 & 35389440 & 0 & 1 \\
F12\_S08 & 168 & high & 1.000000 & 35389440 & 0 & 1 \\
F12\_S09 & 168 & high & 1.000000 & 35389440 & 0 & 1 \\
F12\_S10 & 168 & high & 1.000000 & 35389440 & 0 & 1 \\
F12\_S11 & 168 & high & 1.000000 & 35389440 & 0 & 1 \\
F12\_S12 & 168 & high & 1.000000 & 35389440 & 0 & 1 \\
F12\_S01 & 300 & low & 0.244444 & 9209088 & 87382 & 4 \\
F12\_S02 & 300 & low & 0.244444 & 9209088 & 87382 & 4 \\
F12\_S03 & 300 & low & 0.244444 & 9209088 & 87382 & 4 \\
F12\_S04 & 300 & low & 0.244444 & 9209088 & 87382 & 4 \\
F12\_S05 & 300 & low & 0.244444 & 9209088 & 87382 & 4 \\
F12\_S06 & 300 & low & 0.244444 & 9209088 & 87382 & 4 \\
F12\_S07 & 300 & low & 0.244444 & 9209088 & 87382 & 4 \\
F12\_S08 & 300 & low & 0.244444 & 9209088 & 87382 & 4 \\
F12\_S09 & 300 & low & 0.244444 & 9209088 & 87382 & 4 \\
F12\_S10 & 300 & low & 0.244444 & 9209088 & 87382 & 4 \\
F12\_S11 & 300 & low & 0.244444 & 9209088 & 87382 & 4 \\
F12\_S12 & 300 & low & 0.244444 & 9209088 & 87382 & 4 \\
F12\_S01 & 300 & mid & 0.500000 & 17729280 & 43691 & 2 \\
F12\_S02 & 300 & mid & 0.500000 & 17729280 & 43691 & 2 \\
F12\_S03 & 300 & mid & 0.500000 & 17729280 & 43691 & 2 \\
F12\_S04 & 300 & mid & 0.500000 & 17729280 & 43691 & 2 \\
F12\_S05 & 300 & mid & 0.500000 & 17729280 & 43691 & 2 \\
F12\_S06 & 300 & mid & 0.500000 & 17729280 & 43691 & 2 \\
F12\_S07 & 300 & mid & 0.500000 & 17729280 & 43691 & 2 \\
F12\_S08 & 300 & mid & 0.500000 & 17729280 & 43691 & 2 \\
F12\_S09 & 300 & mid & 0.500000 & 17729280 & 43691 & 2 \\
F12\_S10 & 300 & mid & 0.500000 & 17729280 & 43691 & 2 \\
F12\_S11 & 300 & mid & 0.500000 & 17729280 & 43691 & 2 \\
F12\_S12 & 300 & mid & 0.500000 & 17729280 & 43691 & 2 \\
F12\_S01 & 300 & high & 1.000000 & 35389440 & 0 & 1 \\
F12\_S02 & 300 & high & 1.000000 & 35389440 & 0 & 1 \\
F12\_S03 & 300 & high & 1.000000 & 35389440 & 0 & 1 \\
F12\_S04 & 300 & high & 1.000000 & 35389440 & 0 & 1 \\
F12\_S05 & 300 & high & 1.000000 & 35389440 & 0 & 1 \\
F12\_S06 & 300 & high & 1.000000 & 35389440 & 0 & 1 \\
F12\_S07 & 300 & high & 1.000000 & 35389440 & 0 & 1 \\
F12\_S08 & 300 & high & 1.000000 & 35389440 & 0 & 1 \\
F12\_S09 & 300 & high & 1.000000 & 35389440 & 0 & 1 \\
F12\_S10 & 300 & high & 1.000000 & 35389440 & 0 & 1 \\
F12\_S11 & 300 & high & 1.000000 & 35389440 & 0 & 1 \\
F12\_S12 & 300 & high & 1.000000 & 35389440 & 0 & 1 \\
F12\_S01 & 1000 & low & 0.288889 & 10582272 & 87896 & 7 \\
F12\_S02 & 1000 & low & 0.288889 & 10582272 & 87896 & 7 \\
F12\_S03 & 1000 & low & 0.288889 & 10582272 & 87896 & 7 \\
F12\_S04 & 1000 & low & 0.288889 & 10582272 & 87896 & 7 \\
F12\_S05 & 1000 & low & 0.288889 & 10582272 & 87896 & 7 \\
F12\_S06 & 1000 & low & 0.288889 & 10582272 & 87896 & 7 \\
F12\_S07 & 1000 & low & 0.288889 & 10582272 & 87896 & 7 \\
F12\_S08 & 1000 & low & 0.288889 & 10582272 & 87896 & 7 \\
F12\_S09 & 1000 & low & 0.288889 & 10582272 & 87896 & 7 \\
F12\_S10 & 1000 & low & 0.288889 & 10582272 & 87896 & 7 \\
F12\_S11 & 1000 & low & 0.288889 & 10582272 & 87896 & 7 \\
F12\_S12 & 1000 & low & 0.288889 & 10582272 & 87896 & 7 \\
F12\_S01 & 1000 & mid & 0.500000 & 17729280 & 43691 & 2 \\
F12\_S02 & 1000 & mid & 0.500000 & 17729280 & 43691 & 2 \\
F12\_S03 & 1000 & mid & 0.500000 & 17729280 & 43691 & 2 \\
F12\_S04 & 1000 & mid & 0.500000 & 17729280 & 43691 & 2 \\
F12\_S05 & 1000 & mid & 0.500000 & 17729280 & 43691 & 2 \\
F12\_S06 & 1000 & mid & 0.500000 & 17729280 & 43691 & 2 \\
F12\_S07 & 1000 & mid & 0.500000 & 17729280 & 43691 & 2 \\
F12\_S08 & 1000 & mid & 0.500000 & 17729280 & 43691 & 2 \\
F12\_S09 & 1000 & mid & 0.500000 & 17729280 & 43691 & 2 \\
F12\_S10 & 1000 & mid & 0.500000 & 17729280 & 43691 & 2 \\
F12\_S11 & 1000 & mid & 0.500000 & 17729280 & 43691 & 2 \\
F12\_S12 & 1000 & mid & 0.500000 & 17729280 & 43691 & 2 \\
F12\_S01 & 1000 & high & 1.000000 & 35389440 & 0 & 1 \\
F12\_S02 & 1000 & high & 1.000000 & 35389440 & 0 & 1 \\
F12\_S03 & 1000 & high & 1.000000 & 35389440 & 0 & 1 \\
F12\_S04 & 1000 & high & 1.000000 & 35389440 & 0 & 1 \\
F12\_S05 & 1000 & high & 1.000000 & 35389440 & 0 & 1 \\
F12\_S06 & 1000 & high & 1.000000 & 35389440 & 0 & 1 \\
F12\_S07 & 1000 & high & 1.000000 & 35389440 & 0 & 1 \\
F12\_S08 & 1000 & high & 1.000000 & 35389440 & 0 & 1 \\
F12\_S09 & 1000 & high & 1.000000 & 35389440 & 0 & 1 \\
F12\_S10 & 1000 & high & 1.000000 & 35389440 & 0 & 1 \\
F12\_S11 & 1000 & high & 1.000000 & 35389440 & 0 & 1 \\
F12\_S12 & 1000 & high & 1.000000 & 35389440 & 0 & 1 \\
F12\_S01 & 3000 & low & 0.244444 & 9491712 & 87382 & 4 \\
F12\_S02 & 3000 & low & 0.244444 & 9491712 & 87382 & 4 \\
F12\_S03 & 3000 & low & 0.244444 & 9491712 & 87382 & 4 \\
F12\_S04 & 3000 & low & 0.244444 & 9491712 & 87382 & 4 \\
F12\_S05 & 3000 & low & 0.244444 & 9491712 & 87382 & 4 \\
F12\_S06 & 3000 & low & 0.244444 & 9491712 & 87382 & 4 \\
F12\_S07 & 3000 & low & 0.244444 & 9491712 & 87382 & 4 \\
F12\_S08 & 3000 & low & 0.244444 & 9491712 & 87382 & 4 \\
F12\_S09 & 3000 & low & 0.244444 & 9491712 & 87382 & 4 \\
F12\_S10 & 3000 & low & 0.244444 & 9491712 & 87382 & 4 \\
F12\_S11 & 3000 & low & 0.244444 & 9491712 & 87382 & 4 \\
F12\_S12 & 3000 & low & 0.244444 & 9491712 & 87382 & 4 \\
F12\_S01 & 3000 & mid & 0.500000 & 18247680 & 43691 & 2 \\
F12\_S02 & 3000 & mid & 0.500000 & 18247680 & 43691 & 2 \\
F12\_S03 & 3000 & mid & 0.500000 & 18247680 & 43691 & 2 \\
F12\_S04 & 3000 & mid & 0.500000 & 18247680 & 43691 & 2 \\
F12\_S05 & 3000 & mid & 0.500000 & 18247680 & 43691 & 2 \\
F12\_S06 & 3000 & mid & 0.500000 & 18247680 & 43691 & 2 \\
F12\_S07 & 3000 & mid & 0.500000 & 18247680 & 43691 & 2 \\
F12\_S08 & 3000 & mid & 0.500000 & 18247680 & 43691 & 2 \\
F12\_S09 & 3000 & mid & 0.500000 & 18247680 & 43691 & 2 \\
F12\_S10 & 3000 & mid & 0.500000 & 18247680 & 43691 & 2 \\
F12\_S11 & 3000 & mid & 0.500000 & 18247680 & 43691 & 2 \\
F12\_S12 & 3000 & mid & 0.500000 & 18247680 & 43691 & 2 \\
F12\_S01 & 3000 & high & 1.000000 & 35389440 & 0 & 1 \\
F12\_S02 & 3000 & high & 1.000000 & 35389440 & 0 & 1 \\
F12\_S03 & 3000 & high & 1.000000 & 35389440 & 0 & 1 \\
F12\_S04 & 3000 & high & 1.000000 & 35389440 & 0 & 1 \\
F12\_S05 & 3000 & high & 1.000000 & 35389440 & 0 & 1 \\
F12\_S06 & 3000 & high & 1.000000 & 35389440 & 0 & 1 \\
F12\_S07 & 3000 & high & 1.000000 & 35389440 & 0 & 1 \\
F12\_S08 & 3000 & high & 1.000000 & 35389440 & 0 & 1 \\
F12\_S09 & 3000 & high & 1.000000 & 35389440 & 0 & 1 \\
F12\_S10 & 3000 & high & 1.000000 & 35389440 & 0 & 1 \\
F12\_S11 & 3000 & high & 1.000000 & 35389440 & 0 & 1 \\
F12\_S12 & 3000 & high & 1.000000 & 35389440 & 0 & 1 \\
F12\_S01 & 10000 & low & 0.244444 & 9756672 & 87382 & 4 \\
F12\_S02 & 10000 & low & 0.244444 & 9756672 & 87382 & 4 \\
F12\_S03 & 10000 & low & 0.244444 & 9756672 & 87382 & 4 \\
F12\_S04 & 10000 & low & 0.244444 & 9756672 & 87382 & 4 \\
F12\_S05 & 10000 & low & 0.244444 & 9756672 & 87382 & 4 \\
F12\_S06 & 10000 & low & 0.244444 & 9756672 & 87382 & 4 \\
F12\_S07 & 10000 & low & 0.244444 & 9756672 & 87382 & 4 \\
F12\_S08 & 10000 & low & 0.244444 & 9756672 & 87382 & 4 \\
F12\_S09 & 10000 & low & 0.244444 & 9756672 & 87382 & 4 \\
F12\_S10 & 10000 & low & 0.244444 & 9756672 & 87382 & 4 \\
F12\_S11 & 10000 & low & 0.244444 & 9756672 & 87382 & 4 \\
F12\_S12 & 10000 & low & 0.244444 & 9756672 & 87382 & 4 \\
F12\_S01 & 10000 & mid & 0.500000 & 18247680 & 43691 & 2 \\
F12\_S02 & 10000 & mid & 0.500000 & 18247680 & 43691 & 2 \\
F12\_S03 & 10000 & mid & 0.500000 & 18247680 & 43691 & 2 \\
F12\_S04 & 10000 & mid & 0.500000 & 18247680 & 43691 & 2 \\
F12\_S05 & 10000 & mid & 0.500000 & 18247680 & 43691 & 2 \\
F12\_S06 & 10000 & mid & 0.500000 & 18247680 & 43691 & 2 \\
F12\_S07 & 10000 & mid & 0.500000 & 18247680 & 43691 & 2 \\
F12\_S08 & 10000 & mid & 0.500000 & 18247680 & 43691 & 2 \\
F12\_S09 & 10000 & mid & 0.500000 & 18247680 & 43691 & 2 \\
F12\_S10 & 10000 & mid & 0.500000 & 18247680 & 43691 & 2 \\
F12\_S11 & 10000 & mid & 0.500000 & 18247680 & 43691 & 2 \\
F12\_S12 & 10000 & mid & 0.500000 & 18247680 & 43691 & 2 \\
F12\_S01 & 10000 & high & 1.000000 & 35389440 & 0 & 1 \\
F12\_S02 & 10000 & high & 1.000000 & 35389440 & 0 & 1 \\
F12\_S03 & 10000 & high & 1.000000 & 35389440 & 0 & 1 \\
F12\_S04 & 10000 & high & 1.000000 & 35389440 & 0 & 1 \\
F12\_S05 & 10000 & high & 1.000000 & 35389440 & 0 & 1 \\
F12\_S06 & 10000 & high & 1.000000 & 35389440 & 0 & 1 \\
F12\_S07 & 10000 & high & 1.000000 & 35389440 & 0 & 1 \\
F12\_S08 & 10000 & high & 1.000000 & 35389440 & 0 & 1 \\
F12\_S09 & 10000 & high & 1.000000 & 35389440 & 0 & 1 \\
F12\_S10 & 10000 & high & 1.000000 & 35389440 & 0 & 1 \\
F12\_S11 & 10000 & high & 1.000000 & 35389440 & 0 & 1 \\
F12\_S12 & 10000 & high & 1.000000 & 35389440 & 0 & 1 \\
F12\_S01 & 30000 & low & 0.288889 & 11489280 & 99672 & 7 \\
F12\_S02 & 30000 & low & 0.288889 & 11489280 & 99672 & 7 \\
F12\_S03 & 30000 & low & 0.288889 & 11489280 & 99672 & 7 \\
F12\_S04 & 30000 & low & 0.288889 & 11489280 & 99672 & 7 \\
F12\_S05 & 30000 & low & 0.288889 & 11489280 & 99672 & 7 \\
F12\_S06 & 30000 & low & 0.288889 & 11489280 & 99672 & 7 \\
F12\_S07 & 30000 & low & 0.288889 & 11489280 & 99672 & 7 \\
F12\_S08 & 30000 & low & 0.288889 & 11489280 & 99672 & 7 \\
F12\_S09 & 30000 & low & 0.288889 & 11489280 & 99672 & 7 \\
F12\_S10 & 30000 & low & 0.288889 & 11489280 & 99672 & 7 \\
F12\_S11 & 30000 & low & 0.288889 & 11489280 & 99672 & 7 \\
F12\_S12 & 30000 & low & 0.288889 & 11489280 & 99672 & 7 \\
F12\_S01 & 30000 & mid & 0.500000 & 18800640 & 43691 & 2 \\
F12\_S02 & 30000 & mid & 0.500000 & 18800640 & 43691 & 2 \\
F12\_S03 & 30000 & mid & 0.500000 & 18800640 & 43691 & 2 \\
F12\_S04 & 30000 & mid & 0.500000 & 18800640 & 43691 & 2 \\
F12\_S05 & 30000 & mid & 0.500000 & 18800640 & 43691 & 2 \\
F12\_S06 & 30000 & mid & 0.500000 & 18800640 & 43691 & 2 \\
F12\_S07 & 30000 & mid & 0.500000 & 18800640 & 43691 & 2 \\
F12\_S08 & 30000 & mid & 0.500000 & 18800640 & 43691 & 2 \\
F12\_S09 & 30000 & mid & 0.500000 & 18800640 & 43691 & 2 \\
F12\_S10 & 30000 & mid & 0.500000 & 18800640 & 43691 & 2 \\
F12\_S11 & 30000 & mid & 0.500000 & 18800640 & 43691 & 2 \\
F12\_S12 & 30000 & mid & 0.500000 & 18800640 & 43691 & 2 \\
F12\_S01 & 30000 & high & 1.000000 & 35389440 & 0 & 1 \\
F12\_S02 & 30000 & high & 1.000000 & 35389440 & 0 & 1 \\
F12\_S03 & 30000 & high & 1.000000 & 35389440 & 0 & 1 \\
F12\_S04 & 30000 & high & 1.000000 & 35389440 & 0 & 1 \\
F12\_S05 & 30000 & high & 1.000000 & 35389440 & 0 & 1 \\
F12\_S06 & 30000 & high & 1.000000 & 35389440 & 0 & 1 \\
F12\_S07 & 30000 & high & 1.000000 & 35389440 & 0 & 1 \\
F12\_S08 & 30000 & high & 1.000000 & 35389440 & 0 & 1 \\
F12\_S09 & 30000 & high & 1.000000 & 35389440 & 0 & 1 \\
F12\_S10 & 30000 & high & 1.000000 & 35389440 & 0 & 1 \\
F12\_S11 & 30000 & high & 1.000000 & 35389440 & 0 & 1 \\
F12\_S12 & 30000 & high & 1.000000 & 35389440 & 0 & 1 \\
F12\_S01 & 100000 & low & 0.288889 & 11796480 & 142052 & -- \\
F12\_S02 & 100000 & low & 0.288889 & 11796480 & 142052 & -- \\
F12\_S03 & 100000 & low & 0.288889 & 11796480 & 142052 & -- \\
F12\_S04 & 100000 & low & 0.288889 & 11796480 & 142052 & -- \\
F12\_S05 & 100000 & low & 0.288889 & 11796480 & 142052 & -- \\
F12\_S06 & 100000 & low & 0.288889 & 11796480 & 142052 & -- \\
F12\_S07 & 100000 & low & 0.288889 & 11796480 & 142052 & -- \\
F12\_S08 & 100000 & low & 0.288889 & 11796480 & 142052 & -- \\
F12\_S09 & 100000 & low & 0.288889 & 11796480 & 142052 & -- \\
F12\_S10 & 100000 & low & 0.288889 & 11796480 & 142052 & -- \\
F12\_S11 & 100000 & low & 0.288889 & 11796480 & 142052 & -- \\
F12\_S12 & 100000 & low & 0.288889 & 11796480 & 142052 & -- \\
F12\_S01 & 100000 & mid & 0.622222 & 22855680 & 95575 & 8 \\
F12\_S02 & 100000 & mid & 0.622222 & 22855680 & 95575 & 8 \\
F12\_S03 & 100000 & mid & 0.622222 & 22855680 & 95575 & 8 \\
F12\_S04 & 100000 & mid & 0.622222 & 22855680 & 95575 & 8 \\
F12\_S05 & 100000 & mid & 0.622222 & 22855680 & 95575 & 8 \\
F12\_S06 & 100000 & mid & 0.622222 & 22855680 & 95575 & 8 \\
F12\_S07 & 100000 & mid & 0.622222 & 22855680 & 95575 & 8 \\
F12\_S08 & 100000 & mid & 0.622222 & 22855680 & 95575 & 8 \\
F12\_S09 & 100000 & mid & 0.622222 & 22855680 & 95575 & 8 \\
F12\_S10 & 100000 & mid & 0.622222 & 22855680 & 95575 & 8 \\
F12\_S11 & 100000 & mid & 0.622222 & 22855680 & 95575 & 8 \\
F12\_S12 & 100000 & mid & 0.622222 & 22855680 & 95575 & 8 \\
F12\_S01 & 100000 & high & 1.000000 & 35389440 & 0 & 1 \\
F12\_S02 & 100000 & high & 1.000000 & 35389440 & 0 & 1 \\
F12\_S03 & 100000 & high & 1.000000 & 35389440 & 0 & 1 \\
F12\_S04 & 100000 & high & 1.000000 & 35389440 & 0 & 1 \\
F12\_S05 & 100000 & high & 1.000000 & 35389440 & 0 & 1 \\
F12\_S06 & 100000 & high & 1.000000 & 35389440 & 0 & 1 \\
F12\_S07 & 100000 & high & 1.000000 & 35389440 & 0 & 1 \\
F12\_S08 & 100000 & high & 1.000000 & 35389440 & 0 & 1 \\
F12\_S09 & 100000 & high & 1.000000 & 35389440 & 0 & 1 \\
F12\_S10 & 100000 & high & 1.000000 & 35389440 & 0 & 1 \\
F12\_S11 & 100000 & high & 1.000000 & 35389440 & 0 & 1 \\
F12\_S12 & 100000 & high & 1.000000 & 35389440 & 0 & 1 \\
\bottomrule
\end{longtable}
\end{landscape}

\begin{landscape}
\section{Complete Desynchronization Sensitivity}
Probe phases are independently assigned from fixed seeds. Probes exchange neither action nor queue state, and no joint admission is introduced.
\scriptsize\begin{longtable}{llp{38mm}rrrp{30mm}}
\caption{Complete desynchronization software-sensitivity results (240 retained rows).}\label{tab:desyncfull}\\
\toprule
Seed & Cap. & Method & Core frac. & On-time & Max Q (B) & Phases (ms) \\
\midrule
\endfirsthead
\toprule
Seed & Cap. & Method & Core frac. & On-time & Max Q (B) & Phases (ms) \\
\midrule
\endhead
2026090901 & low & desynchronized\_independent & 0.333333 & 0.000000 & 131029 & 595/193/691 \\
2026090901 & low & synchronized\_independent & 0.244444 & 0.000000 & 131029 & 0/0/0 \\
2026090901 & low & central\_coordinated & 0.000000 & -- & 16341 & 0/0/0 \\
2026090901 & low & shielded\_coordinated & 0.000000 & -- & 16341 & 0/0/0 \\
2026090901 & mid & desynchronized\_independent & 0.666667 & 0.505556 & 130985 & 595/193/691 \\
2026090901 & mid & synchronized\_independent & 0.500000 & 0.000000 & 130985 & 0/0/0 \\
2026090901 & mid & central\_coordinated & 0.333333 & 1.000000 & 130985 & 0/0/0 \\
2026090901 & mid & shielded\_coordinated & 0.333333 & 1.000000 & 130985 & 0/0/0 \\
2026090901 & high & desynchronized\_independent & 1.000000 & 1.000000 & 121111 & 595/193/691 \\
2026090901 & high & synchronized\_independent & 1.000000 & 1.000000 & 130941 & 0/0/0 \\
2026090901 & high & central\_coordinated & 1.000000 & 1.000000 & 130941 & 0/0/0 \\
2026090901 & high & shielded\_coordinated & 1.000000 & 1.000000 & 130941 & 0/0/0 \\
2026090902 & low & desynchronized\_independent & 0.333333 & 0.000000 & 131029 & 403/632/306 \\
2026090902 & low & synchronized\_independent & 0.244444 & 0.000000 & 131029 & 0/0/0 \\
2026090902 & low & central\_coordinated & 0.000000 & -- & 16341 & 0/0/0 \\
2026090902 & low & shielded\_coordinated & 0.000000 & -- & 16341 & 0/0/0 \\
2026090902 & mid & desynchronized\_independent & 0.666667 & 0.500000 & 130985 & 403/632/306 \\
2026090902 & mid & synchronized\_independent & 0.500000 & 0.000000 & 130985 & 0/0/0 \\
2026090902 & mid & central\_coordinated & 0.333333 & 1.000000 & 130985 & 0/0/0 \\
2026090902 & mid & shielded\_coordinated & 0.333333 & 1.000000 & 130985 & 0/0/0 \\
2026090902 & high & desynchronized\_independent & 1.000000 & 1.000000 & 130876 & 403/632/306 \\
2026090902 & high & synchronized\_independent & 1.000000 & 1.000000 & 130941 & 0/0/0 \\
2026090902 & high & central\_coordinated & 1.000000 & 1.000000 & 130941 & 0/0/0 \\
2026090902 & high & shielded\_coordinated & 1.000000 & 1.000000 & 130941 & 0/0/0 \\
2026090903 & low & desynchronized\_independent & 0.333333 & 0.000000 & 131029 & 489/690/996 \\
2026090903 & low & synchronized\_independent & 0.244444 & 0.000000 & 131029 & 0/0/0 \\
2026090903 & low & central\_coordinated & 0.000000 & -- & 16341 & 0/0/0 \\
2026090903 & low & shielded\_coordinated & 0.000000 & -- & 16341 & 0/0/0 \\
2026090903 & mid & desynchronized\_independent & 0.666667 & 0.500000 & 130985 & 489/690/996 \\
2026090903 & mid & synchronized\_independent & 0.500000 & 0.000000 & 130985 & 0/0/0 \\
2026090903 & mid & central\_coordinated & 0.333333 & 1.000000 & 130985 & 0/0/0 \\
2026090903 & mid & shielded\_coordinated & 0.333333 & 1.000000 & 130985 & 0/0/0 \\
2026090903 & high & desynchronized\_independent & 1.000000 & 1.000000 & 1180 & 489/690/996 \\
2026090903 & high & synchronized\_independent & 1.000000 & 1.000000 & 130941 & 0/0/0 \\
2026090903 & high & central\_coordinated & 1.000000 & 1.000000 & 130941 & 0/0/0 \\
2026090903 & high & shielded\_coordinated & 1.000000 & 1.000000 & 130941 & 0/0/0 \\
2026090904 & low & desynchronized\_independent & 0.333333 & 0.000000 & 131028 & 945/680/799 \\
2026090904 & low & synchronized\_independent & 0.244444 & 0.000000 & 131029 & 0/0/0 \\
2026090904 & low & central\_coordinated & 0.000000 & -- & 16341 & 0/0/0 \\
2026090904 & low & shielded\_coordinated & 0.000000 & -- & 16341 & 0/0/0 \\
2026090904 & mid & desynchronized\_independent & 0.666667 & 0.500000 & 130985 & 945/680/799 \\
2026090904 & mid & synchronized\_independent & 0.500000 & 0.000000 & 130985 & 0/0/0 \\
2026090904 & mid & central\_coordinated & 0.333333 & 1.000000 & 130985 & 0/0/0 \\
2026090904 & mid & shielded\_coordinated & 0.333333 & 1.000000 & 130985 & 0/0/0 \\
2026090904 & high & desynchronized\_independent & 1.000000 & 1.000000 & 21234 & 945/680/799 \\
2026090904 & high & synchronized\_independent & 1.000000 & 1.000000 & 130941 & 0/0/0 \\
2026090904 & high & central\_coordinated & 1.000000 & 1.000000 & 130941 & 0/0/0 \\
2026090904 & high & shielded\_coordinated & 1.000000 & 1.000000 & 130941 & 0/0/0 \\
2026090905 & low & desynchronized\_independent & 0.333333 & 0.000000 & 131028 & 154/172/828 \\
2026090905 & low & synchronized\_independent & 0.244444 & 0.000000 & 131029 & 0/0/0 \\
2026090905 & low & central\_coordinated & 0.000000 & -- & 16341 & 0/0/0 \\
2026090905 & low & shielded\_coordinated & 0.000000 & -- & 16341 & 0/0/0 \\
2026090905 & mid & desynchronized\_independent & 0.666667 & 0.505556 & 130985 & 154/172/828 \\
2026090905 & mid & synchronized\_independent & 0.500000 & 0.000000 & 130985 & 0/0/0 \\
2026090905 & mid & central\_coordinated & 0.333333 & 1.000000 & 130985 & 0/0/0 \\
2026090905 & mid & shielded\_coordinated & 0.333333 & 1.000000 & 130985 & 0/0/0 \\
2026090905 & high & desynchronized\_independent & 1.000000 & 1.000000 & 130876 & 154/172/828 \\
2026090905 & high & synchronized\_independent & 1.000000 & 1.000000 & 130941 & 0/0/0 \\
2026090905 & high & central\_coordinated & 1.000000 & 1.000000 & 130941 & 0/0/0 \\
2026090905 & high & shielded\_coordinated & 1.000000 & 1.000000 & 130941 & 0/0/0 \\
2026090906 & low & desynchronized\_independent & 0.333333 & 0.000000 & 131029 & 60/626/317 \\
2026090906 & low & synchronized\_independent & 0.244444 & 0.000000 & 131029 & 0/0/0 \\
2026090906 & low & central\_coordinated & 0.000000 & -- & 16341 & 0/0/0 \\
2026090906 & low & shielded\_coordinated & 0.000000 & -- & 16341 & 0/0/0 \\
2026090906 & mid & desynchronized\_independent & 0.666667 & 0.505556 & 130985 & 60/626/317 \\
2026090906 & mid & synchronized\_independent & 0.500000 & 0.000000 & 130985 & 0/0/0 \\
2026090906 & mid & central\_coordinated & 0.333333 & 1.000000 & 130985 & 0/0/0 \\
2026090906 & mid & shielded\_coordinated & 0.333333 & 1.000000 & 130985 & 0/0/0 \\
2026090906 & high & desynchronized\_independent & 1.000000 & 1.000000 & 130876 & 60/626/317 \\
2026090906 & high & synchronized\_independent & 1.000000 & 1.000000 & 130941 & 0/0/0 \\
2026090906 & high & central\_coordinated & 1.000000 & 1.000000 & 130941 & 0/0/0 \\
2026090906 & high & shielded\_coordinated & 1.000000 & 1.000000 & 130941 & 0/0/0 \\
2026090907 & low & desynchronized\_independent & 0.333333 & 0.000000 & 131029 & 753/955/576 \\
2026090907 & low & synchronized\_independent & 0.244444 & 0.000000 & 131029 & 0/0/0 \\
2026090907 & low & central\_coordinated & 0.000000 & -- & 16341 & 0/0/0 \\
2026090907 & low & shielded\_coordinated & 0.000000 & -- & 16341 & 0/0/0 \\
2026090907 & mid & desynchronized\_independent & 0.666667 & 0.500000 & 130985 & 753/955/576 \\
2026090907 & mid & synchronized\_independent & 0.500000 & 0.000000 & 130985 & 0/0/0 \\
2026090907 & mid & central\_coordinated & 0.333333 & 1.000000 & 130985 & 0/0/0 \\
2026090907 & mid & shielded\_coordinated & 0.333333 & 1.000000 & 130985 & 0/0/0 \\
2026090907 & high & desynchronized\_independent & 1.000000 & 1.000000 & 17302 & 753/955/576 \\
2026090907 & high & synchronized\_independent & 1.000000 & 1.000000 & 130941 & 0/0/0 \\
2026090907 & high & central\_coordinated & 1.000000 & 1.000000 & 130941 & 0/0/0 \\
2026090907 & high & shielded\_coordinated & 1.000000 & 1.000000 & 130941 & 0/0/0 \\
2026090908 & low & desynchronized\_independent & 0.333333 & 0.000000 & 131029 & 884/483/389 \\
2026090908 & low & synchronized\_independent & 0.244444 & 0.000000 & 131029 & 0/0/0 \\
2026090908 & low & central\_coordinated & 0.000000 & -- & 16341 & 0/0/0 \\
2026090908 & low & shielded\_coordinated & 0.000000 & -- & 16341 & 0/0/0 \\
2026090908 & mid & desynchronized\_independent & 0.666667 & 0.500000 & 130985 & 884/483/389 \\
2026090908 & mid & synchronized\_independent & 0.500000 & 0.000000 & 130985 & 0/0/0 \\
2026090908 & mid & central\_coordinated & 0.333333 & 1.000000 & 130985 & 0/0/0 \\
2026090908 & mid & shielded\_coordinated & 0.333333 & 1.000000 & 130985 & 0/0/0 \\
2026090908 & high & desynchronized\_independent & 1.000000 & 1.000000 & 71828 & 884/483/389 \\
2026090908 & high & synchronized\_independent & 1.000000 & 1.000000 & 130941 & 0/0/0 \\
2026090908 & high & central\_coordinated & 1.000000 & 1.000000 & 130941 & 0/0/0 \\
2026090908 & high & shielded\_coordinated & 1.000000 & 1.000000 & 130941 & 0/0/0 \\
2026090909 & low & desynchronized\_independent & 0.333333 & 0.000000 & 131029 & 530/348/834 \\
2026090909 & low & synchronized\_independent & 0.244444 & 0.000000 & 131029 & 0/0/0 \\
2026090909 & low & central\_coordinated & 0.000000 & -- & 16341 & 0/0/0 \\
2026090909 & low & shielded\_coordinated & 0.000000 & -- & 16341 & 0/0/0 \\
2026090909 & mid & desynchronized\_independent & 0.666667 & 0.500000 & 130985 & 530/348/834 \\
2026090909 & mid & synchronized\_independent & 0.500000 & 0.000000 & 130985 & 0/0/0 \\
2026090909 & mid & central\_coordinated & 0.333333 & 1.000000 & 130985 & 0/0/0 \\
2026090909 & mid & shielded\_coordinated & 0.333333 & 1.000000 & 130985 & 0/0/0 \\
2026090909 & high & desynchronized\_independent & 1.000000 & 1.000000 & 64881 & 530/348/834 \\
2026090909 & high & synchronized\_independent & 1.000000 & 1.000000 & 130941 & 0/0/0 \\
2026090909 & high & central\_coordinated & 1.000000 & 1.000000 & 130941 & 0/0/0 \\
2026090909 & high & shielded\_coordinated & 1.000000 & 1.000000 & 130941 & 0/0/0 \\
2026090910 & low & desynchronized\_independent & 0.333333 & 0.000000 & 131028 & 611/131/96 \\
2026090910 & low & synchronized\_independent & 0.244444 & 0.000000 & 131029 & 0/0/0 \\
2026090910 & low & central\_coordinated & 0.000000 & -- & 16341 & 0/0/0 \\
2026090910 & low & shielded\_coordinated & 0.000000 & -- & 16341 & 0/0/0 \\
2026090910 & mid & desynchronized\_independent & 0.666667 & 0.505556 & 130985 & 611/131/96 \\
2026090910 & mid & synchronized\_independent & 0.500000 & 0.000000 & 130985 & 0/0/0 \\
2026090910 & mid & central\_coordinated & 0.333333 & 1.000000 & 130985 & 0/0/0 \\
2026090910 & mid & shielded\_coordinated & 0.333333 & 1.000000 & 130985 & 0/0/0 \\
2026090910 & high & desynchronized\_independent & 1.000000 & 1.000000 & 130876 & 611/131/96 \\
2026090910 & high & synchronized\_independent & 1.000000 & 1.000000 & 130941 & 0/0/0 \\
2026090910 & high & central\_coordinated & 1.000000 & 1.000000 & 130941 & 0/0/0 \\
2026090910 & high & shielded\_coordinated & 1.000000 & 1.000000 & 130941 & 0/0/0 \\
2026090911 & low & desynchronized\_independent & 0.333333 & 0.000000 & 131029 & 286/697/613 \\
2026090911 & low & synchronized\_independent & 0.244444 & 0.000000 & 131029 & 0/0/0 \\
2026090911 & low & central\_coordinated & 0.000000 & -- & 16341 & 0/0/0 \\
2026090911 & low & shielded\_coordinated & 0.000000 & -- & 16341 & 0/0/0 \\
2026090911 & mid & desynchronized\_independent & 0.666667 & 0.505556 & 130985 & 286/697/613 \\
2026090911 & mid & synchronized\_independent & 0.500000 & 0.000000 & 130985 & 0/0/0 \\
2026090911 & mid & central\_coordinated & 0.333333 & 1.000000 & 130985 & 0/0/0 \\
2026090911 & mid & shielded\_coordinated & 0.333333 & 1.000000 & 130985 & 0/0/0 \\
2026090911 & high & desynchronized\_independent & 1.000000 & 1.000000 & 118752 & 286/697/613 \\
2026090911 & high & synchronized\_independent & 1.000000 & 1.000000 & 130941 & 0/0/0 \\
2026090911 & high & central\_coordinated & 1.000000 & 1.000000 & 130941 & 0/0/0 \\
2026090911 & high & shielded\_coordinated & 1.000000 & 1.000000 & 130941 & 0/0/0 \\
2026090912 & low & desynchronized\_independent & 0.333333 & 0.000000 & 131028 & 72/898/820 \\
2026090912 & low & synchronized\_independent & 0.244444 & 0.000000 & 131029 & 0/0/0 \\
2026090912 & low & central\_coordinated & 0.000000 & -- & 16341 & 0/0/0 \\
2026090912 & low & shielded\_coordinated & 0.000000 & -- & 16341 & 0/0/0 \\
2026090912 & mid & desynchronized\_independent & 0.666667 & 0.505556 & 130985 & 72/898/820 \\
2026090912 & mid & synchronized\_independent & 0.500000 & 0.000000 & 130985 & 0/0/0 \\
2026090912 & mid & central\_coordinated & 0.333333 & 1.000000 & 130985 & 0/0/0 \\
2026090912 & mid & shielded\_coordinated & 0.333333 & 1.000000 & 130985 & 0/0/0 \\
2026090912 & high & desynchronized\_independent & 1.000000 & 1.000000 & 102368 & 72/898/820 \\
2026090912 & high & synchronized\_independent & 1.000000 & 1.000000 & 130941 & 0/0/0 \\
2026090912 & high & central\_coordinated & 1.000000 & 1.000000 & 130941 & 0/0/0 \\
2026090912 & high & shielded\_coordinated & 1.000000 & 1.000000 & 130941 & 0/0/0 \\
2026090913 & low & desynchronized\_independent & 0.333333 & 0.000000 & 131028 & 31/997/848 \\
2026090913 & low & synchronized\_independent & 0.244444 & 0.000000 & 131029 & 0/0/0 \\
2026090913 & low & central\_coordinated & 0.000000 & -- & 16341 & 0/0/0 \\
2026090913 & low & shielded\_coordinated & 0.000000 & -- & 16341 & 0/0/0 \\
2026090913 & mid & desynchronized\_independent & 0.666667 & 0.505556 & 130985 & 31/997/848 \\
2026090913 & mid & synchronized\_independent & 0.500000 & 0.000000 & 130985 & 0/0/0 \\
2026090913 & mid & central\_coordinated & 0.333333 & 1.000000 & 130985 & 0/0/0 \\
2026090913 & mid & shielded\_coordinated & 0.333333 & 1.000000 & 130985 & 0/0/0 \\
2026090913 & high & desynchronized\_independent & 1.000000 & 1.000000 & 118490 & 31/997/848 \\
2026090913 & high & synchronized\_independent & 1.000000 & 1.000000 & 130941 & 0/0/0 \\
2026090913 & high & central\_coordinated & 1.000000 & 1.000000 & 130941 & 0/0/0 \\
2026090913 & high & shielded\_coordinated & 1.000000 & 1.000000 & 130941 & 0/0/0 \\
2026090914 & low & desynchronized\_independent & 0.333333 & 0.000000 & 131029 & 379/396/678 \\
2026090914 & low & synchronized\_independent & 0.244444 & 0.000000 & 131029 & 0/0/0 \\
2026090914 & low & central\_coordinated & 0.000000 & -- & 16341 & 0/0/0 \\
2026090914 & low & shielded\_coordinated & 0.000000 & -- & 16341 & 0/0/0 \\
2026090914 & mid & desynchronized\_independent & 0.666667 & 0.500000 & 130985 & 379/396/678 \\
2026090914 & mid & synchronized\_independent & 0.500000 & 0.000000 & 130985 & 0/0/0 \\
2026090914 & mid & central\_coordinated & 0.333333 & 1.000000 & 130985 & 0/0/0 \\
2026090914 & mid & shielded\_coordinated & 0.333333 & 1.000000 & 130985 & 0/0/0 \\
2026090914 & high & desynchronized\_independent & 1.000000 & 1.000000 & 126223 & 379/396/678 \\
2026090914 & high & synchronized\_independent & 1.000000 & 1.000000 & 130941 & 0/0/0 \\
2026090914 & high & central\_coordinated & 1.000000 & 1.000000 & 130941 & 0/0/0 \\
2026090914 & high & shielded\_coordinated & 1.000000 & 1.000000 & 130941 & 0/0/0 \\
2026090915 & low & desynchronized\_independent & 0.333333 & 0.000000 & 131028 & 212/604/213 \\
2026090915 & low & synchronized\_independent & 0.244444 & 0.000000 & 131029 & 0/0/0 \\
2026090915 & low & central\_coordinated & 0.000000 & -- & 16341 & 0/0/0 \\
2026090915 & low & shielded\_coordinated & 0.000000 & -- & 16341 & 0/0/0 \\
2026090915 & mid & desynchronized\_independent & 0.666667 & 0.505556 & 130985 & 212/604/213 \\
2026090915 & mid & synchronized\_independent & 0.500000 & 0.000000 & 130985 & 0/0/0 \\
2026090915 & mid & central\_coordinated & 0.333333 & 1.000000 & 130985 & 0/0/0 \\
2026090915 & mid & shielded\_coordinated & 0.333333 & 1.000000 & 130985 & 0/0/0 \\
2026090915 & high & desynchronized\_independent & 1.000000 & 1.000000 & 130876 & 212/604/213 \\
2026090915 & high & synchronized\_independent & 1.000000 & 1.000000 & 130941 & 0/0/0 \\
2026090915 & high & central\_coordinated & 1.000000 & 1.000000 & 130941 & 0/0/0 \\
2026090915 & high & shielded\_coordinated & 1.000000 & 1.000000 & 130941 & 0/0/0 \\
2026090916 & low & desynchronized\_independent & 0.333333 & 0.000000 & 131029 & 586/124/90 \\
2026090916 & low & synchronized\_independent & 0.244444 & 0.000000 & 131029 & 0/0/0 \\
2026090916 & low & central\_coordinated & 0.000000 & -- & 16341 & 0/0/0 \\
2026090916 & low & shielded\_coordinated & 0.000000 & -- & 16341 & 0/0/0 \\
2026090916 & mid & desynchronized\_independent & 0.666667 & 0.505556 & 130985 & 586/124/90 \\
2026090916 & mid & synchronized\_independent & 0.500000 & 0.000000 & 130985 & 0/0/0 \\
2026090916 & mid & central\_coordinated & 0.333333 & 1.000000 & 130985 & 0/0/0 \\
2026090916 & mid & shielded\_coordinated & 0.333333 & 1.000000 & 130985 & 0/0/0 \\
2026090916 & high & desynchronized\_independent & 1.000000 & 1.000000 & 130876 & 586/124/90 \\
2026090916 & high & synchronized\_independent & 1.000000 & 1.000000 & 130941 & 0/0/0 \\
2026090916 & high & central\_coordinated & 1.000000 & 1.000000 & 130941 & 0/0/0 \\
2026090916 & high & shielded\_coordinated & 1.000000 & 1.000000 & 130941 & 0/0/0 \\
2026090917 & low & desynchronized\_independent & 0.333333 & 0.000000 & 131029 & 13/5/654 \\
2026090917 & low & synchronized\_independent & 0.244444 & 0.000000 & 131029 & 0/0/0 \\
2026090917 & low & central\_coordinated & 0.000000 & -- & 16341 & 0/0/0 \\
2026090917 & low & shielded\_coordinated & 0.000000 & -- & 16341 & 0/0/0 \\
2026090917 & mid & desynchronized\_independent & 0.666667 & 0.500000 & 130985 & 13/5/654 \\
2026090917 & mid & synchronized\_independent & 0.500000 & 0.000000 & 130985 & 0/0/0 \\
2026090917 & mid & central\_coordinated & 0.333333 & 1.000000 & 130985 & 0/0/0 \\
2026090917 & mid & shielded\_coordinated & 0.333333 & 1.000000 & 130985 & 0/0/0 \\
2026090917 & high & desynchronized\_independent & 1.000000 & 1.000000 & 130876 & 13/5/654 \\
2026090917 & high & synchronized\_independent & 1.000000 & 1.000000 & 130941 & 0/0/0 \\
2026090917 & high & central\_coordinated & 1.000000 & 1.000000 & 130941 & 0/0/0 \\
2026090917 & high & shielded\_coordinated & 1.000000 & 1.000000 & 130941 & 0/0/0 \\
2026090918 & low & desynchronized\_independent & 0.333333 & 0.000000 & 131029 & 303/966/432 \\
2026090918 & low & synchronized\_independent & 0.244444 & 0.000000 & 131029 & 0/0/0 \\
2026090918 & low & central\_coordinated & 0.000000 & -- & 16341 & 0/0/0 \\
2026090918 & low & shielded\_coordinated & 0.000000 & -- & 16341 & 0/0/0 \\
2026090918 & mid & desynchronized\_independent & 0.666667 & 0.500000 & 130985 & 303/966/432 \\
2026090918 & mid & synchronized\_independent & 0.500000 & 0.000000 & 130985 & 0/0/0 \\
2026090918 & mid & central\_coordinated & 0.333333 & 1.000000 & 130985 & 0/0/0 \\
2026090918 & mid & shielded\_coordinated & 0.333333 & 1.000000 & 130985 & 0/0/0 \\
2026090918 & high & desynchronized\_independent & 1.000000 & 1.000000 & 91882 & 303/966/432 \\
2026090918 & high & synchronized\_independent & 1.000000 & 1.000000 & 130941 & 0/0/0 \\
2026090918 & high & central\_coordinated & 1.000000 & 1.000000 & 130941 & 0/0/0 \\
2026090918 & high & shielded\_coordinated & 1.000000 & 1.000000 & 130941 & 0/0/0 \\
2026090919 & low & desynchronized\_independent & 0.333333 & 0.000000 & 131029 & 94/427/937 \\
2026090919 & low & synchronized\_independent & 0.244444 & 0.000000 & 131029 & 0/0/0 \\
2026090919 & low & central\_coordinated & 0.000000 & -- & 16341 & 0/0/0 \\
2026090919 & low & shielded\_coordinated & 0.000000 & -- & 16341 & 0/0/0 \\
2026090919 & mid & desynchronized\_independent & 0.666667 & 0.500000 & 130985 & 94/427/937 \\
2026090919 & mid & synchronized\_independent & 0.500000 & 0.000000 & 130985 & 0/0/0 \\
2026090919 & mid & central\_coordinated & 0.333333 & 1.000000 & 130985 & 0/0/0 \\
2026090919 & mid & shielded\_coordinated & 0.333333 & 1.000000 & 130985 & 0/0/0 \\
2026090919 & high & desynchronized\_independent & 1.000000 & 1.000000 & 93848 & 94/427/937 \\
2026090919 & high & synchronized\_independent & 1.000000 & 1.000000 & 130941 & 0/0/0 \\
2026090919 & high & central\_coordinated & 1.000000 & 1.000000 & 130941 & 0/0/0 \\
2026090919 & high & shielded\_coordinated & 1.000000 & 1.000000 & 130941 & 0/0/0 \\
2026090920 & low & desynchronized\_independent & 0.333333 & 0.000000 & 131029 & 699/148/926 \\
2026090920 & low & synchronized\_independent & 0.244444 & 0.000000 & 131029 & 0/0/0 \\
2026090920 & low & central\_coordinated & 0.000000 & -- & 16341 & 0/0/0 \\
2026090920 & low & shielded\_coordinated & 0.000000 & -- & 16341 & 0/0/0 \\
2026090920 & mid & desynchronized\_independent & 0.666667 & 0.505556 & 130985 & 699/148/926 \\
2026090920 & mid & synchronized\_independent & 0.500000 & 0.000000 & 130985 & 0/0/0 \\
2026090920 & mid & central\_coordinated & 0.333333 & 1.000000 & 130985 & 0/0/0 \\
2026090920 & mid & shielded\_coordinated & 0.333333 & 1.000000 & 130985 & 0/0/0 \\
2026090920 & high & desynchronized\_independent & 1.000000 & 1.000000 & 72483 & 699/148/926 \\
2026090920 & high & synchronized\_independent & 1.000000 & 1.000000 & 130941 & 0/0/0 \\
2026090920 & high & central\_coordinated & 1.000000 & 1.000000 & 130941 & 0/0/0 \\
2026090920 & high & shielded\_coordinated & 1.000000 & 1.000000 & 130941 & 0/0/0 \\
\bottomrule
\end{longtable}
\end{landscape}

\begin{landscape}
\section{Complete Nominal-Deadline Reconstruction}
Post-release latency is result return minus actual event creation. Nominal latency is result return minus scheduled release. The 1-s quantity in the main paper is a post-release service SLO, not a complete schedule-to-result deadline.
\tiny\begin{longtable}{llllrrrrrr}
\caption{Complete accepted-hardware timing reconstruction (648 retained cells; latency in ms).}\label{tab:timingfull}\\
\toprule
Session & Method & Cap. & Dev. & $n_c$ & Post & Nom. & Late p50 & p95 & p99 \\
\midrule
\endfirsthead
\toprule
Session & Method & Cap. & Dev. & $n_c$ & Post & Nom. & Late p50 & p95 & p99 \\
\midrule
\endhead
F12\_S01 & fixed\_local & low & P1 & 0 & -- & -- & 1862.2 & 1920.9 & 1924.1 \\
F12\_S01 & fixed\_local & low & P2 & 0 & -- & -- & 760.3 & 831.8 & 834.1 \\
F12\_S01 & fixed\_local & low & P3 & 0 & -- & -- & 86.3 & 155.7 & 162.0 \\
F12\_S01 & fixed\_local & mid & P1 & 0 & -- & -- & 1828.4 & 1979.1 & 1980.6 \\
F12\_S01 & fixed\_local & mid & P2 & 0 & -- & -- & 636.8 & 775.7 & 780.4 \\
F12\_S01 & fixed\_local & mid & P3 & 0 & -- & -- & 87.5 & 165.1 & 172.5 \\
F12\_S01 & fixed\_local & high & P1 & 0 & -- & -- & 1827.5 & 1863.0 & 1913.4 \\
F12\_S01 & fixed\_local & high & P2 & 0 & -- & -- & 637.0 & 671.0 & 673.3 \\
F12\_S01 & fixed\_local & high & P3 & 0 & -- & -- & 84.4 & 156.5 & 163.5 \\
F12\_S01 & fixed\_summary & low & P1 & 0 & -- & -- & 1952.6 & 2002.0 & 2006.5 \\
F12\_S01 & fixed\_summary & low & P2 & 0 & -- & -- & 712.1 & 750.1 & 753.4 \\
F12\_S01 & fixed\_summary & low & P3 & 0 & -- & -- & 84.0 & 156.4 & 163.8 \\
F12\_S01 & fixed\_summary & mid & P1 & 0 & -- & -- & 1922.6 & 1997.1 & 2001.5 \\
F12\_S01 & fixed\_summary & mid & P2 & 0 & -- & -- & 716.0 & 804.2 & 806.5 \\
F12\_S01 & fixed\_summary & mid & P3 & 0 & -- & -- & 85.2 & 153.6 & 160.0 \\
F12\_S01 & fixed\_summary & high & P1 & 0 & -- & -- & 1925.6 & 1978.0 & 1981.4 \\
F12\_S01 & fixed\_summary & high & P2 & 0 & -- & -- & 745.5 & 802.9 & 806.8 \\
F12\_S01 & fixed\_summary & high & P3 & 0 & -- & -- & 88.3 & 159.6 & 166.3 \\
F12\_S01 & fixed\_full & low & P1 & 90 & 0.0000 & 0.0000 & 1887.7 & 1981.0 & 1983.3 \\
F12\_S01 & fixed\_full & low & P2 & 90 & 0.0000 & 0.0000 & 736.5 & 827.9 & 830.4 \\
F12\_S01 & fixed\_full & low & P3 & 90 & 0.0000 & 0.0000 & 79.3 & 139.5 & 145.3 \\
F12\_S01 & fixed\_full & mid & P1 & 90 & 0.0000 & 0.0000 & 1859.1 & 2483.1 & 2486.7 \\
F12\_S01 & fixed\_full & mid & P2 & 90 & 0.0000 & 0.0000 & 686.7 & 1341.8 & 1346.6 \\
F12\_S01 & fixed\_full & mid & P3 & 90 & 0.0000 & 0.0000 & 79.2 & 141.9 & 148.5 \\
F12\_S01 & fixed\_full & high & P1 & 90 & 0.9667 & 0.0000 & 2024.3 & 2101.8 & 2105.2 \\
F12\_S01 & fixed\_full & high & P2 & 90 & 0.9222 & 0.0000 & 830.5 & 948.8 & 952.3 \\
F12\_S01 & fixed\_full & high & P3 & 90 & 1.0000 & 1.0000 & 78.6 & 141.2 & 145.6 \\
F12\_S01 & independent\_dpp & low & P1 & 22 & 0.0000 & 0.0000 & 1851.5 & 2183.8 & 2301.2 \\
F12\_S01 & independent\_dpp & low & P2 & 22 & 0.0000 & 0.0000 & 661.8 & 792.4 & 1053.9 \\
F12\_S01 & independent\_dpp & low & P3 & 22 & 0.0000 & 0.0000 & 82.3 & 152.1 & 156.3 \\
F12\_S01 & independent\_dpp & mid & P1 & 45 & 0.0000 & 0.0000 & 1998.0 & 2108.8 & 2111.0 \\
F12\_S01 & independent\_dpp & mid & P2 & 45 & 0.0000 & 0.0000 & 754.9 & 822.0 & 824.8 \\
F12\_S01 & independent\_dpp & mid & P3 & 45 & 0.0000 & 0.0000 & 81.4 & 146.9 & 152.4 \\
F12\_S01 & independent\_dpp & high & P1 & 90 & 0.9444 & 0.0000 & 1967.9 & 2006.6 & 2009.6 \\
F12\_S01 & independent\_dpp & high & P2 & 90 & 0.9556 & 0.0000 & 797.6 & 836.1 & 841.8 \\
F12\_S01 & independent\_dpp & high & P3 & 90 & 1.0000 & 1.0000 & 78.4 & 136.1 & 140.8 \\
F12\_S01 & central\_dpp & low & P1 & 0 & -- & -- & 2049.7 & 2151.4 & 2153.5 \\
F12\_S01 & central\_dpp & low & P2 & 0 & -- & -- & 858.4 & 942.6 & 947.1 \\
F12\_S01 & central\_dpp & low & P3 & 0 & -- & -- & 83.2 & 148.0 & 154.0 \\
F12\_S01 & central\_dpp & mid & P1 & 30 & 0.8000 & 0.0000 & 1910.4 & 1992.9 & 1995.4 \\
F12\_S01 & central\_dpp & mid & P2 & 30 & 0.8000 & 0.0000 & 732.4 & 809.6 & 813.0 \\
F12\_S01 & central\_dpp & mid & P3 & 30 & 1.0000 & 1.0000 & 83.1 & 152.5 & 159.3 \\
F12\_S01 & central\_dpp & high & P1 & 90 & 1.0000 & 0.0000 & 2070.5 & 2105.4 & 2109.1 \\
F12\_S01 & central\_dpp & high & P2 & 90 & 0.9889 & 0.0000 & 907.5 & 949.5 & 953.2 \\
F12\_S01 & central\_dpp & high & P3 & 90 & 1.0000 & 1.0000 & 78.5 & 144.9 & 150.8 \\
F12\_S01 & shielded\_dpp & low & P1 & 0 & -- & -- & 1771.7 & 1832.4 & 1835.5 \\
F12\_S01 & shielded\_dpp & low & P2 & 0 & -- & -- & 678.1 & 749.6 & 786.5 \\
F12\_S01 & shielded\_dpp & low & P3 & 0 & -- & -- & 87.8 & 156.7 & 162.3 \\
F12\_S01 & shielded\_dpp & mid & P1 & 30 & 0.9000 & 0.0000 & 1934.5 & 2160.6 & 2163.7 \\
F12\_S01 & shielded\_dpp & mid & P2 & 30 & 1.0000 & 0.0000 & 856.4 & 964.9 & 967.9 \\
F12\_S01 & shielded\_dpp & mid & P3 & 30 & 1.0000 & 1.0000 & 86.3 & 149.8 & 155.9 \\
F12\_S01 & shielded\_dpp & high & P1 & 90 & 0.9333 & 0.0000 & 1979.2 & 2065.8 & 2070.3 \\
F12\_S01 & shielded\_dpp & high & P2 & 90 & 0.9444 & 0.0000 & 876.4 & 913.1 & 974.2 \\
F12\_S01 & shielded\_dpp & high & P3 & 90 & 1.0000 & 1.0000 & 79.7 & 136.0 & 140.8 \\
F12\_S02 & fixed\_local & low & P1 & 0 & -- & -- & 1886.4 & 2021.6 & 2024.1 \\
F12\_S02 & fixed\_local & low & P2 & 0 & -- & -- & 594.0 & 740.2 & 743.2 \\
F12\_S02 & fixed\_local & low & P3 & 0 & -- & -- & 88.2 & 164.6 & 171.5 \\
F12\_S02 & fixed\_local & mid & P1 & 0 & -- & -- & 2003.6 & 2046.7 & 2051.4 \\
F12\_S02 & fixed\_local & mid & P2 & 0 & -- & -- & 717.1 & 761.0 & 764.1 \\
F12\_S02 & fixed\_local & mid & P3 & 0 & -- & -- & 92.2 & 161.7 & 168.4 \\
F12\_S02 & fixed\_local & high & P1 & 0 & -- & -- & 1896.5 & 1938.2 & 1941.3 \\
F12\_S02 & fixed\_local & high & P2 & 0 & -- & -- & 737.1 & 775.7 & 778.9 \\
F12\_S02 & fixed\_local & high & P3 & 0 & -- & -- & 91.3 & 160.5 & 167.6 \\
F12\_S02 & fixed\_summary & low & P1 & 0 & -- & -- & 1725.1 & 1764.9 & 1767.8 \\
F12\_S02 & fixed\_summary & low & P2 & 0 & -- & -- & 550.6 & 585.7 & 589.6 \\
F12\_S02 & fixed\_summary & low & P3 & 0 & -- & -- & 90.6 & 160.8 & 167.3 \\
F12\_S02 & fixed\_summary & mid & P1 & 0 & -- & -- & 1915.3 & 2065.5 & 2072.1 \\
F12\_S02 & fixed\_summary & mid & P2 & 0 & -- & -- & 677.3 & 1103.5 & 1106.6 \\
F12\_S02 & fixed\_summary & mid & P3 & 0 & -- & -- & 85.8 & 156.6 & 162.7 \\
F12\_S02 & fixed\_summary & high & P1 & 0 & -- & -- & 1878.3 & 1949.3 & 1959.8 \\
F12\_S02 & fixed\_summary & high & P2 & 0 & -- & -- & 713.1 & 787.0 & 798.8 \\
F12\_S02 & fixed\_summary & high & P3 & 0 & -- & -- & 89.0 & 158.6 & 165.8 \\
F12\_S02 & fixed\_full & low & P1 & 90 & 0.0000 & 0.0000 & 1929.0 & 2110.4 & 2114.6 \\
F12\_S02 & fixed\_full & low & P2 & 90 & 0.0000 & 0.0000 & 763.0 & 876.5 & 879.5 \\
F12\_S02 & fixed\_full & low & P3 & 90 & 0.0000 & 0.0000 & 78.2 & 135.9 & 142.5 \\
F12\_S02 & fixed\_full & mid & P1 & 90 & 0.0000 & 0.0000 & 2251.6 & 2305.8 & 2309.9 \\
F12\_S02 & fixed\_full & mid & P2 & 90 & 0.0000 & 0.0000 & 616.5 & 733.8 & 750.5 \\
F12\_S02 & fixed\_full & mid & P3 & 90 & 0.0000 & 0.0000 & 76.4 & 136.6 & 142.5 \\
F12\_S02 & fixed\_full & high & P1 & 90 & 1.0000 & 0.0000 & 1969.5 & 1996.6 & 2001.5 \\
F12\_S02 & fixed\_full & high & P2 & 90 & 0.9889 & 0.0000 & 827.2 & 865.1 & 867.9 \\
F12\_S02 & fixed\_full & high & P3 & 90 & 1.0000 & 1.0000 & 80.0 & 145.5 & 151.7 \\
F12\_S02 & independent\_dpp & low & P1 & 22 & 0.0000 & 0.0000 & 1906.4 & 2060.6 & 2064.0 \\
F12\_S02 & independent\_dpp & low & P2 & 22 & 0.0000 & 0.0000 & 748.0 & 899.5 & 904.3 \\
F12\_S02 & independent\_dpp & low & P3 & 22 & 0.0000 & 0.0000 & 78.9 & 142.0 & 147.5 \\
F12\_S02 & independent\_dpp & mid & P1 & 45 & 0.0000 & 0.0000 & 1980.3 & 2013.9 & 2015.6 \\
F12\_S02 & independent\_dpp & mid & P2 & 45 & 0.0000 & 0.0000 & 703.8 & 732.1 & 734.0 \\
F12\_S02 & independent\_dpp & mid & P3 & 45 & 0.0000 & 0.0000 & 79.5 & 147.8 & 152.5 \\
F12\_S02 & independent\_dpp & high & P1 & 90 & 0.9111 & 0.0000 & 1982.0 & 2052.8 & 2056.0 \\
F12\_S02 & independent\_dpp & high & P2 & 90 & 0.9333 & 0.0000 & 787.2 & 831.0 & 832.9 \\
F12\_S02 & independent\_dpp & high & P3 & 90 & 1.0000 & 1.0000 & 78.0 & 137.5 & 141.7 \\
F12\_S02 & central\_dpp & low & P1 & 0 & -- & -- & 2017.6 & 2154.2 & 2157.5 \\
F12\_S02 & central\_dpp & low & P2 & 0 & -- & -- & 738.5 & 861.7 & 865.2 \\
F12\_S02 & central\_dpp & low & P3 & 0 & -- & -- & 87.9 & 158.3 & 164.1 \\
F12\_S02 & central\_dpp & mid & P1 & 30 & 0.7333 & 0.0000 & 1936.7 & 2071.6 & 2074.8 \\
F12\_S02 & central\_dpp & mid & P2 & 30 & 0.9333 & 0.0000 & 881.1 & 913.6 & 916.2 \\
F12\_S02 & central\_dpp & mid & P3 & 30 & 1.0000 & 1.0000 & 86.2 & 150.4 & 156.4 \\
F12\_S02 & central\_dpp & high & P1 & 90 & 0.9222 & 0.0000 & 1942.0 & 1995.4 & 2000.0 \\
F12\_S02 & central\_dpp & high & P2 & 90 & 0.9444 & 0.0000 & 816.7 & 862.3 & 865.6 \\
F12\_S02 & central\_dpp & high & P3 & 90 & 1.0000 & 1.0000 & 81.3 & 145.8 & 151.8 \\
F12\_S02 & shielded\_dpp & low & P1 & 0 & -- & -- & 1931.4 & 1995.7 & 1999.1 \\
F12\_S02 & shielded\_dpp & low & P2 & 0 & -- & -- & 842.4 & 903.9 & 908.0 \\
F12\_S02 & shielded\_dpp & low & P3 & 0 & -- & -- & 86.6 & 160.6 & 168.3 \\
F12\_S02 & shielded\_dpp & mid & P1 & 30 & 0.9000 & 0.0000 & 2120.6 & 2150.6 & 2151.7 \\
F12\_S02 & shielded\_dpp & mid & P2 & 30 & 0.9000 & 0.0000 & 793.9 & 826.4 & 828.9 \\
F12\_S02 & shielded\_dpp & mid & P3 & 30 & 1.0000 & 1.0000 & 80.3 & 148.0 & 155.9 \\
F12\_S02 & shielded\_dpp & high & P1 & 90 & 0.9444 & 0.0000 & 1954.9 & 2007.4 & 2039.4 \\
F12\_S02 & shielded\_dpp & high & P2 & 90 & 0.9333 & 0.0000 & 781.0 & 858.6 & 873.6 \\
F12\_S02 & shielded\_dpp & high & P3 & 90 & 1.0000 & 1.0000 & 81.5 & 140.4 & 145.8 \\
F12\_S03 & fixed\_local & low & P1 & 0 & -- & -- & 1892.3 & 1942.8 & 1950.8 \\
F12\_S03 & fixed\_local & low & P2 & 0 & -- & -- & 676.8 & 732.8 & 744.0 \\
F12\_S03 & fixed\_local & low & P3 & 0 & -- & -- & 90.0 & 162.6 & 169.3 \\
F12\_S03 & fixed\_local & mid & P1 & 0 & -- & -- & 1823.3 & 1868.5 & 1871.1 \\
F12\_S03 & fixed\_local & mid & P2 & 0 & -- & -- & 736.6 & 837.7 & 840.2 \\
F12\_S03 & fixed\_local & mid & P3 & 0 & -- & -- & 89.2 & 160.3 & 166.4 \\
F12\_S03 & fixed\_local & high & P1 & 0 & -- & -- & 1939.5 & 2052.7 & 2061.3 \\
F12\_S03 & fixed\_local & high & P2 & 0 & -- & -- & 724.4 & 838.4 & 852.7 \\
F12\_S03 & fixed\_local & high & P3 & 0 & -- & -- & 90.8 & 167.3 & 174.3 \\
F12\_S03 & fixed\_summary & low & P1 & 0 & -- & -- & 1926.4 & 1967.6 & 1971.1 \\
F12\_S03 & fixed\_summary & low & P2 & 0 & -- & -- & 710.5 & 757.2 & 762.6 \\
F12\_S03 & fixed\_summary & low & P3 & 0 & -- & -- & 90.1 & 166.0 & 172.4 \\
F12\_S03 & fixed\_summary & mid & P1 & 0 & -- & -- & 2190.5 & 2234.0 & 2237.6 \\
F12\_S03 & fixed\_summary & mid & P2 & 0 & -- & -- & 877.1 & 918.5 & 922.6 \\
F12\_S03 & fixed\_summary & mid & P3 & 0 & -- & -- & 90.8 & 162.8 & 169.3 \\
F12\_S03 & fixed\_summary & high & P1 & 0 & -- & -- & 2024.1 & 2062.5 & 2066.0 \\
F12\_S03 & fixed\_summary & high & P2 & 0 & -- & -- & 748.8 & 787.7 & 790.7 \\
F12\_S03 & fixed\_summary & high & P3 & 0 & -- & -- & 89.3 & 160.2 & 166.3 \\
F12\_S03 & fixed\_full & low & P1 & 90 & 0.0000 & 0.0000 & 1963.8 & 2094.5 & 2097.1 \\
F12\_S03 & fixed\_full & low & P2 & 90 & 0.0000 & 0.0000 & 864.3 & 926.3 & 947.9 \\
F12\_S03 & fixed\_full & low & P3 & 90 & 0.0000 & 0.0000 & 78.2 & 147.4 & 153.9 \\
F12\_S03 & fixed\_full & mid & P1 & 90 & 0.0000 & 0.0000 & 1903.6 & 1975.4 & 1976.5 \\
F12\_S03 & fixed\_full & mid & P2 & 90 & 0.0000 & 0.0000 & 752.0 & 824.8 & 829.8 \\
F12\_S03 & fixed\_full & mid & P3 & 90 & 0.0000 & 0.0000 & 80.7 & 141.3 & 147.6 \\
F12\_S03 & fixed\_full & high & P1 & 90 & 0.9667 & 0.0000 & 2006.1 & 2050.0 & 2098.4 \\
F12\_S03 & fixed\_full & high & P2 & 90 & 0.9333 & 0.0000 & 847.8 & 892.2 & 1004.3 \\
F12\_S03 & fixed\_full & high & P3 & 90 & 1.0000 & 1.0000 & 86.9 & 150.9 & 156.0 \\
F12\_S03 & independent\_dpp & low & P1 & 22 & 0.0000 & 0.0000 & 1937.5 & 2023.7 & 2025.5 \\
F12\_S03 & independent\_dpp & low & P2 & 22 & 0.0000 & 0.0000 & 755.3 & 832.0 & 841.6 \\
F12\_S03 & independent\_dpp & low & P3 & 22 & 0.0000 & 0.0000 & 82.0 & 151.6 & 158.6 \\
F12\_S03 & independent\_dpp & mid & P1 & 45 & 0.0000 & 0.0000 & 7335.7 & 7390.7 & 7395.4 \\
F12\_S03 & independent\_dpp & mid & P2 & 45 & 0.0000 & 0.0000 & 6236.0 & 6290.3 & 6295.7 \\
F12\_S03 & independent\_dpp & mid & P3 & 45 & 0.0000 & 0.0000 & 4851.8 & 4918.9 & 4924.2 \\
F12\_S03 & independent\_dpp & high & P1 & 90 & 0.9889 & 0.0000 & 1867.0 & 2014.7 & 2018.8 \\
F12\_S03 & independent\_dpp & high & P2 & 90 & 0.9556 & 0.0000 & 699.5 & 814.5 & 817.0 \\
F12\_S03 & independent\_dpp & high & P3 & 90 & 1.0000 & 1.0000 & 74.3 & 134.1 & 139.7 \\
F12\_S03 & central\_dpp & low & P1 & 0 & -- & -- & 1984.5 & 2024.7 & 2030.2 \\
F12\_S03 & central\_dpp & low & P2 & 0 & -- & -- & 838.4 & 875.1 & 878.3 \\
F12\_S03 & central\_dpp & low & P3 & 0 & -- & -- & 90.1 & 161.4 & 168.3 \\
F12\_S03 & central\_dpp & mid & P1 & 30 & 0.8667 & 0.0000 & 1949.1 & 2028.9 & 2031.1 \\
F12\_S03 & central\_dpp & mid & P2 & 30 & 0.7667 & 0.0000 & 915.3 & 954.2 & 957.1 \\
F12\_S03 & central\_dpp & mid & P3 & 30 & 1.0000 & 1.0000 & 81.0 & 147.8 & 154.9 \\
F12\_S03 & central\_dpp & high & P1 & 90 & 0.9667 & 0.0000 & 1963.1 & 2023.1 & 2025.9 \\
F12\_S03 & central\_dpp & high & P2 & 90 & 1.0000 & 0.0000 & 856.1 & 892.5 & 896.6 \\
F12\_S03 & central\_dpp & high & P3 & 90 & 1.0000 & 1.0000 & 73.2 & 133.8 & 140.1 \\
F12\_S03 & shielded\_dpp & low & P1 & 0 & -- & -- & 2074.6 & 2116.7 & 2120.3 \\
F12\_S03 & shielded\_dpp & low & P2 & 0 & -- & -- & 627.9 & 665.5 & 668.1 \\
F12\_S03 & shielded\_dpp & low & P3 & 0 & -- & -- & 84.1 & 154.8 & 160.5 \\
F12\_S03 & shielded\_dpp & mid & P1 & 30 & 0.8000 & 0.0000 & 2003.3 & 2147.9 & 2150.1 \\
F12\_S03 & shielded\_dpp & mid & P2 & 30 & 0.8333 & 0.0000 & 758.1 & 862.1 & 892.8 \\
F12\_S03 & shielded\_dpp & mid & P3 & 30 & 1.0000 & 1.0000 & 86.6 & 151.2 & 156.3 \\
F12\_S03 & shielded\_dpp & high & P1 & 90 & 0.9333 & 0.0000 & 2057.5 & 2130.4 & 2139.2 \\
F12\_S03 & shielded\_dpp & high & P2 & 90 & 0.9778 & 0.0000 & 939.6 & 989.1 & 992.6 \\
F12\_S03 & shielded\_dpp & high & P3 & 90 & 1.0000 & 1.0000 & 77.9 & 142.0 & 146.9 \\
F12\_S04 & fixed\_local & low & P1 & 0 & -- & -- & 1879.6 & 1975.4 & 1977.9 \\
F12\_S04 & fixed\_local & low & P2 & 0 & -- & -- & 657.9 & 744.4 & 746.8 \\
F12\_S04 & fixed\_local & low & P3 & 0 & -- & -- & 91.4 & 163.7 & 170.4 \\
F12\_S04 & fixed\_local & mid & P1 & 0 & -- & -- & 1944.5 & 2017.8 & 2022.3 \\
F12\_S04 & fixed\_local & mid & P2 & 0 & -- & -- & 724.9 & 795.4 & 798.7 \\
F12\_S04 & fixed\_local & mid & P3 & 0 & -- & -- & 95.3 & 169.4 & 176.3 \\
F12\_S04 & fixed\_local & high & P1 & 0 & -- & -- & 1896.6 & 1934.6 & 1938.1 \\
F12\_S04 & fixed\_local & high & P2 & 0 & -- & -- & 743.8 & 780.2 & 783.8 \\
F12\_S04 & fixed\_local & high & P3 & 0 & -- & -- & 87.9 & 160.4 & 166.0 \\
F12\_S04 & fixed\_summary & low & P1 & 0 & -- & -- & 1958.8 & 2049.4 & 2059.2 \\
F12\_S04 & fixed\_summary & low & P2 & 0 & -- & -- & 719.0 & 800.6 & 806.0 \\
F12\_S04 & fixed\_summary & low & P3 & 0 & -- & -- & 87.6 & 156.2 & 162.2 \\
F12\_S04 & fixed\_summary & mid & P1 & 0 & -- & -- & 1929.8 & 2022.4 & 2027.0 \\
F12\_S04 & fixed\_summary & mid & P2 & 0 & -- & -- & 706.0 & 830.5 & 834.7 \\
F12\_S04 & fixed\_summary & mid & P3 & 0 & -- & -- & 89.1 & 158.3 & 164.2 \\
F12\_S04 & fixed\_summary & high & P1 & 0 & -- & -- & 1973.7 & 2015.9 & 2019.2 \\
F12\_S04 & fixed\_summary & high & P2 & 0 & -- & -- & 595.9 & 629.4 & 630.5 \\
F12\_S04 & fixed\_summary & high & P3 & 0 & -- & -- & 88.2 & 155.6 & 161.3 \\
F12\_S04 & fixed\_full & low & P1 & 90 & 0.0000 & 0.0000 & 2135.3 & 2299.8 & 2301.8 \\
F12\_S04 & fixed\_full & low & P2 & 90 & 0.0000 & 0.0000 & 843.4 & 1033.6 & 1036.2 \\
F12\_S04 & fixed\_full & low & P3 & 90 & 0.0000 & 0.0000 & 78.4 & 140.7 & 146.6 \\
F12\_S04 & fixed\_full & mid & P1 & 90 & 0.0000 & 0.0000 & 1990.7 & 2057.3 & 2060.3 \\
F12\_S04 & fixed\_full & mid & P2 & 90 & 0.0000 & 0.0000 & 864.7 & 897.3 & 900.8 \\
F12\_S04 & fixed\_full & mid & P3 & 90 & 0.0000 & 0.0000 & 78.5 & 141.6 & 146.8 \\
F12\_S04 & fixed\_full & high & P1 & 90 & 0.9556 & 0.0000 & 1968.2 & 2168.1 & 2172.8 \\
F12\_S04 & fixed\_full & high & P2 & 90 & 0.9222 & 0.0000 & 746.4 & 946.3 & 950.5 \\
F12\_S04 & fixed\_full & high & P3 & 90 & 1.0000 & 1.0000 & 78.6 & 142.2 & 146.8 \\
F12\_S04 & independent\_dpp & low & P1 & 22 & 0.0000 & 0.0000 & 2113.4 & 2196.6 & 2199.2 \\
F12\_S04 & independent\_dpp & low & P2 & 22 & 0.0000 & 0.0000 & 1004.2 & 1036.6 & 1038.8 \\
F12\_S04 & independent\_dpp & low & P3 & 22 & 0.0000 & 0.0000 & 85.1 & 154.4 & 159.8 \\
F12\_S04 & independent\_dpp & mid & P1 & 45 & 0.0000 & 0.0000 & 1990.7 & 2228.1 & 2232.6 \\
F12\_S04 & independent\_dpp & mid & P2 & 45 & 0.0000 & 0.0000 & 837.2 & 975.7 & 977.9 \\
F12\_S04 & independent\_dpp & mid & P3 & 45 & 0.0000 & 0.0000 & 77.9 & 143.7 & 149.7 \\
F12\_S04 & independent\_dpp & high & P1 & 90 & 0.9444 & 0.0000 & 2007.0 & 2091.5 & 2094.1 \\
F12\_S04 & independent\_dpp & high & P2 & 90 & 0.9556 & 0.0000 & 806.0 & 927.9 & 929.9 \\
F12\_S04 & independent\_dpp & high & P3 & 90 & 1.0000 & 1.0000 & 80.4 & 145.8 & 151.8 \\
F12\_S04 & central\_dpp & low & P1 & 0 & -- & -- & 1985.7 & 2026.6 & 2029.4 \\
F12\_S04 & central\_dpp & low & P2 & 0 & -- & -- & 669.0 & 706.4 & 709.7 \\
F12\_S04 & central\_dpp & low & P3 & 0 & -- & -- & 85.8 & 158.2 & 163.9 \\
F12\_S04 & central\_dpp & mid & P1 & 30 & 0.6667 & 0.0000 & 1908.9 & 1971.8 & 1982.7 \\
F12\_S04 & central\_dpp & mid & P2 & 30 & 0.7667 & 0.0000 & 747.3 & 842.3 & 846.1 \\
F12\_S04 & central\_dpp & mid & P3 & 30 & 1.0000 & 1.0000 & 86.9 & 156.1 & 162.2 \\
F12\_S04 & central\_dpp & high & P1 & 90 & 0.9556 & 0.0000 & 2089.1 & 2217.9 & 2220.9 \\
F12\_S04 & central\_dpp & high & P2 & 90 & 0.9889 & 0.0000 & 855.7 & 922.1 & 923.9 \\
F12\_S04 & central\_dpp & high & P3 & 90 & 1.0000 & 1.0000 & 74.0 & 133.8 & 139.0 \\
F12\_S04 & shielded\_dpp & low & P1 & 0 & -- & -- & 1873.0 & 1929.6 & 1934.2 \\
F12\_S04 & shielded\_dpp & low & P2 & 0 & -- & -- & 648.1 & 713.7 & 716.7 \\
F12\_S04 & shielded\_dpp & low & P3 & 0 & -- & -- & 85.0 & 156.3 & 162.7 \\
F12\_S04 & shielded\_dpp & mid & P1 & 30 & 0.7667 & 0.0000 & 1978.4 & 2044.0 & 2047.9 \\
F12\_S04 & shielded\_dpp & mid & P2 & 30 & 0.7667 & 0.0000 & 795.4 & 884.6 & 887.1 \\
F12\_S04 & shielded\_dpp & mid & P3 & 30 & 1.0000 & 1.0000 & 81.0 & 144.6 & 150.5 \\
F12\_S04 & shielded\_dpp & high & P1 & 90 & 0.9778 & 0.0000 & 1987.3 & 2025.7 & 2029.0 \\
F12\_S04 & shielded\_dpp & high & P2 & 90 & 0.9111 & 0.0000 & 724.4 & 854.0 & 858.5 \\
F12\_S04 & shielded\_dpp & high & P3 & 90 & 1.0000 & 1.0000 & 78.4 & 144.4 & 150.5 \\
F12\_S05 & fixed\_local & low & P1 & 0 & -- & -- & 2011.8 & 2048.3 & 2050.5 \\
F12\_S05 & fixed\_local & low & P2 & 0 & -- & -- & 735.0 & 769.7 & 771.7 \\
F12\_S05 & fixed\_local & low & P3 & 0 & -- & -- & 88.0 & 158.3 & 164.5 \\
F12\_S05 & fixed\_local & mid & P1 & 0 & -- & -- & 1826.1 & 1896.2 & 1901.4 \\
F12\_S05 & fixed\_local & mid & P2 & 0 & -- & -- & 665.2 & 704.7 & 706.8 \\
F12\_S05 & fixed\_local & mid & P3 & 0 & -- & -- & 86.6 & 158.2 & 164.3 \\
F12\_S05 & fixed\_local & high & P1 & 0 & -- & -- & 1917.4 & 1963.7 & 1966.4 \\
F12\_S05 & fixed\_local & high & P2 & 0 & -- & -- & 763.8 & 800.2 & 802.8 \\
F12\_S05 & fixed\_local & high & P3 & 0 & -- & -- & 90.4 & 169.1 & 175.5 \\
F12\_S05 & fixed\_summary & low & P1 & 0 & -- & -- & 1836.5 & 1895.6 & 1899.3 \\
F12\_S05 & fixed\_summary & low & P2 & 0 & -- & -- & 624.0 & 695.4 & 698.3 \\
F12\_S05 & fixed\_summary & low & P3 & 0 & -- & -- & 88.9 & 160.6 & 166.5 \\
F12\_S05 & fixed\_summary & mid & P1 & 0 & -- & -- & 1787.3 & 1904.6 & 1908.2 \\
F12\_S05 & fixed\_summary & mid & P2 & 0 & -- & -- & 749.4 & 864.3 & 868.8 \\
F12\_S05 & fixed\_summary & mid & P3 & 0 & -- & -- & 91.0 & 163.6 & 170.1 \\
F12\_S05 & fixed\_summary & high & P1 & 0 & -- & -- & 1869.3 & 1931.6 & 1935.0 \\
F12\_S05 & fixed\_summary & high & P2 & 0 & -- & -- & 712.6 & 774.4 & 779.4 \\
F12\_S05 & fixed\_summary & high & P3 & 0 & -- & -- & 86.0 & 152.7 & 159.9 \\
F12\_S05 & fixed\_full & low & P1 & 90 & 0.0000 & 0.0000 & 2017.7 & 2049.2 & 2051.7 \\
F12\_S05 & fixed\_full & low & P2 & 90 & 0.0000 & 0.0000 & 779.2 & 840.4 & 844.0 \\
F12\_S05 & fixed\_full & low & P3 & 90 & 0.0000 & 0.0000 & 76.6 & 140.7 & 146.2 \\
F12\_S05 & fixed\_full & mid & P1 & 90 & 0.0000 & 0.0000 & 1920.6 & 2001.0 & 2049.4 \\
F12\_S05 & fixed\_full & mid & P2 & 90 & 0.0000 & 0.0000 & 746.0 & 797.5 & 843.3 \\
F12\_S05 & fixed\_full & mid & P3 & 90 & 0.0000 & 0.0000 & 78.4 & 144.6 & 151.0 \\
F12\_S05 & fixed\_full & high & P1 & 90 & 0.9111 & 0.0000 & 1861.7 & 1922.9 & 1926.9 \\
F12\_S05 & fixed\_full & high & P2 & 90 & 0.9222 & 0.0000 & 777.4 & 829.5 & 831.5 \\
F12\_S05 & fixed\_full & high & P3 & 90 & 1.0000 & 1.0000 & 81.1 & 140.9 & 146.8 \\
F12\_S05 & independent\_dpp & low & P1 & 22 & 0.0000 & 0.0000 & 1901.3 & 1976.4 & 1981.7 \\
F12\_S05 & independent\_dpp & low & P2 & 22 & 0.0000 & 0.0000 & 654.1 & 712.2 & 715.8 \\
F12\_S05 & independent\_dpp & low & P3 & 22 & 0.0000 & 0.0000 & 81.8 & 145.1 & 151.4 \\
F12\_S05 & independent\_dpp & mid & P1 & 45 & 0.0000 & 0.0000 & 1841.5 & 1921.5 & 1925.4 \\
F12\_S05 & independent\_dpp & mid & P2 & 45 & 0.0000 & 0.0000 & 675.7 & 737.0 & 747.7 \\
F12\_S05 & independent\_dpp & mid & P3 & 45 & 0.0000 & 0.0000 & 74.9 & 137.5 & 142.5 \\
F12\_S05 & independent\_dpp & high & P1 & 90 & 0.9778 & 0.0000 & 1962.0 & 2083.4 & 2087.9 \\
F12\_S05 & independent\_dpp & high & P2 & 90 & 0.9778 & 0.0000 & 886.6 & 920.4 & 923.2 \\
F12\_S05 & independent\_dpp & high & P3 & 90 & 1.0000 & 1.0000 & 76.8 & 139.6 & 145.5 \\
F12\_S05 & central\_dpp & low & P1 & 0 & -- & -- & 1957.7 & 1999.1 & 2002.0 \\
F12\_S05 & central\_dpp & low & P2 & 0 & -- & -- & 769.8 & 811.2 & 814.7 \\
F12\_S05 & central\_dpp & low & P3 & 0 & -- & -- & 85.3 & 153.6 & 160.1 \\
F12\_S05 & central\_dpp & mid & P1 & 30 & 1.0000 & 0.0000 & 1912.7 & 1940.1 & 1942.0 \\
F12\_S05 & central\_dpp & mid & P2 & 30 & 0.8000 & 0.0000 & 793.2 & 835.2 & 837.4 \\
F12\_S05 & central\_dpp & mid & P3 & 30 & 1.0000 & 1.0000 & 82.5 & 150.1 & 155.5 \\
F12\_S05 & central\_dpp & high & P1 & 90 & 0.9111 & 0.0000 & 1946.2 & 1994.5 & 1999.5 \\
F12\_S05 & central\_dpp & high & P2 & 90 & 0.9778 & 0.0000 & 792.1 & 861.7 & 865.5 \\
F12\_S05 & central\_dpp & high & P3 & 90 & 1.0000 & 1.0000 & 77.4 & 141.8 & 146.7 \\
F12\_S05 & shielded\_dpp & low & P1 & 0 & -- & -- & 1893.0 & 1978.6 & 1982.3 \\
F12\_S05 & shielded\_dpp & low & P2 & 0 & -- & -- & 673.0 & 754.2 & 757.7 \\
F12\_S05 & shielded\_dpp & low & P3 & 0 & -- & -- & 84.9 & 152.9 & 159.8 \\
F12\_S05 & shielded\_dpp & mid & P1 & 30 & 0.5667 & 0.0000 & 1903.8 & 1973.6 & 1976.1 \\
F12\_S05 & shielded\_dpp & mid & P2 & 30 & 0.8333 & 0.0000 & 822.2 & 890.1 & 892.5 \\
F12\_S05 & shielded\_dpp & mid & P3 & 30 & 1.0000 & 1.0000 & 81.1 & 146.8 & 152.3 \\
F12\_S05 & shielded\_dpp & high & P1 & 90 & 0.9222 & 0.0000 & 1914.0 & 1990.0 & 1992.6 \\
F12\_S05 & shielded\_dpp & high & P2 & 90 & 0.9667 & 0.0000 & 701.7 & 774.2 & 778.5 \\
F12\_S05 & shielded\_dpp & high & P3 & 90 & 1.0000 & 1.0000 & 80.0 & 140.8 & 147.1 \\
F12\_S06 & fixed\_local & low & P1 & 0 & -- & -- & 1932.5 & 2004.1 & 2017.5 \\
F12\_S06 & fixed\_local & low & P2 & 0 & -- & -- & 752.8 & 824.8 & 837.3 \\
F12\_S06 & fixed\_local & low & P3 & 0 & -- & -- & 90.0 & 164.6 & 171.9 \\
F12\_S06 & fixed\_local & mid & P1 & 0 & -- & -- & 1937.2 & 2042.4 & 2059.8 \\
F12\_S06 & fixed\_local & mid & P2 & 0 & -- & -- & 673.6 & 769.8 & 782.5 \\
F12\_S06 & fixed\_local & mid & P3 & 0 & -- & -- & 90.2 & 159.5 & 166.9 \\
F12\_S06 & fixed\_local & high & P1 & 0 & -- & -- & 1909.3 & 1956.4 & 1959.2 \\
F12\_S06 & fixed\_local & high & P2 & 0 & -- & -- & 785.0 & 826.4 & 829.8 \\
F12\_S06 & fixed\_local & high & P3 & 0 & -- & -- & 90.5 & 162.6 & 169.1 \\
F12\_S06 & fixed\_summary & low & P1 & 0 & -- & -- & 1785.0 & 1821.6 & 1825.8 \\
F12\_S06 & fixed\_summary & low & P2 & 0 & -- & -- & 614.8 & 652.1 & 654.9 \\
F12\_S06 & fixed\_summary & low & P3 & 0 & -- & -- & 88.8 & 157.6 & 164.1 \\
F12\_S06 & fixed\_summary & mid & P1 & 0 & -- & -- & 1889.4 & 1925.6 & 1929.1 \\
F12\_S06 & fixed\_summary & mid & P2 & 0 & -- & -- & 741.5 & 779.2 & 781.7 \\
F12\_S06 & fixed\_summary & mid & P3 & 0 & -- & -- & 85.5 & 156.6 & 163.6 \\
F12\_S06 & fixed\_summary & high & P1 & 0 & -- & -- & 1854.5 & 1988.9 & 2023.9 \\
F12\_S06 & fixed\_summary & high & P2 & 0 & -- & -- & 659.0 & 780.9 & 817.9 \\
F12\_S06 & fixed\_summary & high & P3 & 0 & -- & -- & 91.7 & 165.0 & 171.2 \\
F12\_S06 & fixed\_full & low & P1 & 90 & 0.0000 & 0.0000 & 2001.6 & 2059.9 & 2106.9 \\
F12\_S06 & fixed\_full & low & P2 & 90 & 0.0000 & 0.0000 & 694.5 & 809.4 & 850.0 \\
F12\_S06 & fixed\_full & low & P3 & 90 & 0.0000 & 0.0000 & 80.2 & 143.3 & 148.5 \\
F12\_S06 & fixed\_full & mid & P1 & 90 & 0.0000 & 0.0000 & 1890.5 & 2019.8 & 2023.1 \\
F12\_S06 & fixed\_full & mid & P2 & 90 & 0.0000 & 0.0000 & 652.4 & 715.4 & 735.8 \\
F12\_S06 & fixed\_full & mid & P3 & 90 & 0.0000 & 0.0000 & 77.1 & 141.6 & 146.7 \\
F12\_S06 & fixed\_full & high & P1 & 90 & 0.9667 & 0.0000 & 1966.0 & 2015.7 & 2038.6 \\
F12\_S06 & fixed\_full & high & P2 & 90 & 0.9111 & 0.0000 & 734.7 & 797.3 & 799.7 \\
F12\_S06 & fixed\_full & high & P3 & 90 & 1.0000 & 1.0000 & 80.4 & 144.3 & 150.8 \\
F12\_S06 & independent\_dpp & low & P1 & 22 & 0.0000 & 0.0000 & 2005.8 & 2040.9 & 2043.7 \\
F12\_S06 & independent\_dpp & low & P2 & 22 & 0.0000 & 0.0000 & 696.2 & 797.4 & 800.0 \\
F12\_S06 & independent\_dpp & low & P3 & 22 & 0.0000 & 0.0000 & 84.9 & 156.4 & 161.7 \\
F12\_S06 & independent\_dpp & mid & P1 & 45 & 0.0000 & 0.0000 & 1969.6 & 1996.1 & 2000.6 \\
F12\_S06 & independent\_dpp & mid & P2 & 45 & 0.0000 & 0.0000 & 903.6 & 932.0 & 934.8 \\
F12\_S06 & independent\_dpp & mid & P3 & 45 & 0.0000 & 0.0000 & 79.1 & 144.5 & 150.1 \\
F12\_S06 & independent\_dpp & high & P1 & 90 & 0.9222 & 0.0000 & 1912.8 & 2042.7 & 2046.9 \\
F12\_S06 & independent\_dpp & high & P2 & 90 & 0.9444 & 0.0000 & 730.8 & 856.4 & 859.5 \\
F12\_S06 & independent\_dpp & high & P3 & 90 & 1.0000 & 1.0000 & 73.9 & 132.8 & 138.6 \\
F12\_S06 & central\_dpp & low & P1 & 0 & -- & -- & 1945.7 & 2067.9 & 2070.5 \\
F12\_S06 & central\_dpp & low & P2 & 0 & -- & -- & 746.1 & 825.0 & 827.8 \\
F12\_S06 & central\_dpp & low & P3 & 0 & -- & -- & 84.7 & 151.8 & 158.0 \\
F12\_S06 & central\_dpp & mid & P1 & 30 & 0.9000 & 0.0000 & 2001.6 & 2059.9 & 2070.6 \\
F12\_S06 & central\_dpp & mid & P2 & 30 & 0.8667 & 0.0000 & 869.7 & 935.4 & 937.7 \\
F12\_S06 & central\_dpp & mid & P3 & 30 & 1.0000 & 1.0000 & 89.2 & 151.7 & 158.0 \\
F12\_S06 & central\_dpp & high & P1 & 90 & 0.9444 & 0.0000 & 2175.3 & 2214.4 & 2217.3 \\
F12\_S06 & central\_dpp & high & P2 & 90 & 0.9667 & 0.0000 & 1059.9 & 1102.3 & 1105.7 \\
F12\_S06 & central\_dpp & high & P3 & 90 & 1.0000 & 1.0000 & 79.4 & 154.0 & 160.2 \\
F12\_S06 & shielded\_dpp & low & P1 & 0 & -- & -- & 1933.6 & 1972.0 & 1976.6 \\
F12\_S06 & shielded\_dpp & low & P2 & 0 & -- & -- & 775.1 & 817.3 & 820.9 \\
F12\_S06 & shielded\_dpp & low & P3 & 0 & -- & -- & 89.9 & 160.3 & 166.5 \\
F12\_S06 & shielded\_dpp & mid & P1 & 30 & 0.8333 & 0.0000 & 1830.6 & 1960.8 & 1963.5 \\
F12\_S06 & shielded\_dpp & mid & P2 & 30 & 0.9000 & 0.0000 & 688.1 & 761.8 & 764.0 \\
F12\_S06 & shielded\_dpp & mid & P3 & 30 & 1.0000 & 1.0000 & 82.1 & 150.9 & 156.6 \\
F12\_S06 & shielded\_dpp & high & P1 & 90 & 0.9778 & 0.0000 & 1922.9 & 1970.4 & 1972.8 \\
F12\_S06 & shielded\_dpp & high & P2 & 90 & 0.9556 & 0.0000 & 698.8 & 740.1 & 744.4 \\
F12\_S06 & shielded\_dpp & high & P3 & 90 & 1.0000 & 1.0000 & 88.2 & 149.6 & 155.1 \\
F12\_S07 & fixed\_local & low & P1 & 0 & -- & -- & 1944.9 & 2011.2 & 2014.7 \\
F12\_S07 & fixed\_local & low & P2 & 0 & -- & -- & 730.1 & 809.6 & 812.8 \\
F12\_S07 & fixed\_local & low & P3 & 0 & -- & -- & 88.8 & 160.6 & 167.3 \\
F12\_S07 & fixed\_local & mid & P1 & 0 & -- & -- & 1868.9 & 1943.8 & 1948.2 \\
F12\_S07 & fixed\_local & mid & P2 & 0 & -- & -- & 629.0 & 699.1 & 702.5 \\
F12\_S07 & fixed\_local & mid & P3 & 0 & -- & -- & 93.9 & 164.7 & 171.4 \\
F12\_S07 & fixed\_local & high & P1 & 0 & -- & -- & 2100.5 & 2245.8 & 2249.1 \\
F12\_S07 & fixed\_local & high & P2 & 0 & -- & -- & 895.1 & 1020.1 & 1024.0 \\
F12\_S07 & fixed\_local & high & P3 & 0 & -- & -- & 88.5 & 159.7 & 166.7 \\
F12\_S07 & fixed\_summary & low & P1 & 0 & -- & -- & 1997.9 & 2132.3 & 2136.3 \\
F12\_S07 & fixed\_summary & low & P2 & 0 & -- & -- & 675.1 & 801.3 & 804.5 \\
F12\_S07 & fixed\_summary & low & P3 & 0 & -- & -- & 84.1 & 151.1 & 156.9 \\
F12\_S07 & fixed\_summary & mid & P1 & 0 & -- & -- & 1927.0 & 1978.8 & 1985.2 \\
F12\_S07 & fixed\_summary & mid & P2 & 0 & -- & -- & 612.5 & 680.9 & 684.0 \\
F12\_S07 & fixed\_summary & mid & P3 & 0 & -- & -- & 87.5 & 157.8 & 165.0 \\
F12\_S07 & fixed\_summary & high & P1 & 0 & -- & -- & 2181.8 & 2223.8 & 2227.5 \\
F12\_S07 & fixed\_summary & high & P2 & 0 & -- & -- & 613.1 & 751.0 & 755.5 \\
F12\_S07 & fixed\_summary & high & P3 & 0 & -- & -- & 84.2 & 154.1 & 161.1 \\
F12\_S07 & fixed\_full & low & P1 & 90 & 0.0000 & 0.0000 & 1825.8 & 1945.2 & 1947.3 \\
F12\_S07 & fixed\_full & low & P2 & 90 & 0.0000 & 0.0000 & 557.2 & 701.7 & 705.8 \\
F12\_S07 & fixed\_full & low & P3 & 90 & 0.0000 & 0.0000 & 81.1 & 143.6 & 149.2 \\
F12\_S07 & fixed\_full & mid & P1 & 90 & 0.0000 & 0.0000 & 2023.6 & 2274.3 & 2277.8 \\
F12\_S07 & fixed\_full & mid & P2 & 90 & 0.0000 & 0.0000 & 771.6 & 952.8 & 955.2 \\
F12\_S07 & fixed\_full & mid & P3 & 90 & 0.0000 & 0.0000 & 77.4 & 139.8 & 145.4 \\
F12\_S07 & fixed\_full & high & P1 & 90 & 0.9778 & 0.0000 & 1983.4 & 2025.4 & 2028.5 \\
F12\_S07 & fixed\_full & high & P2 & 90 & 0.9333 & 0.0000 & 793.0 & 850.9 & 872.5 \\
F12\_S07 & fixed\_full & high & P3 & 90 & 1.0000 & 1.0000 & 79.3 & 145.9 & 150.9 \\
F12\_S07 & independent\_dpp & low & P1 & 22 & 0.0000 & 0.0000 & 1909.6 & 2101.8 & 2103.8 \\
F12\_S07 & independent\_dpp & low & P2 & 22 & 0.0000 & 0.0000 & 717.8 & 874.3 & 876.8 \\
F12\_S07 & independent\_dpp & low & P3 & 22 & 0.0000 & 0.0000 & 83.9 & 149.3 & 155.4 \\
F12\_S07 & independent\_dpp & mid & P1 & 45 & 0.0000 & 0.0000 & 2049.3 & 2104.3 & 2107.5 \\
F12\_S07 & independent\_dpp & mid & P2 & 45 & 0.0000 & 0.0000 & 941.2 & 966.4 & 970.3 \\
F12\_S07 & independent\_dpp & mid & P3 & 45 & 0.0000 & 0.0000 & 85.9 & 151.4 & 156.4 \\
F12\_S07 & independent\_dpp & high & P1 & 90 & 0.9667 & 0.0000 & 1967.2 & 2170.5 & 2176.7 \\
F12\_S07 & independent\_dpp & high & P2 & 90 & 0.8889 & 0.0000 & 782.6 & 931.8 & 1045.7 \\
F12\_S07 & independent\_dpp & high & P3 & 90 & 0.9889 & 0.9889 & 81.2 & 328.2 & 333.1 \\
F12\_S07 & central\_dpp & low & P1 & 0 & -- & -- & 2036.6 & 2077.1 & 2082.3 \\
F12\_S07 & central\_dpp & low & P2 & 0 & -- & -- & 723.5 & 764.1 & 768.0 \\
F12\_S07 & central\_dpp & low & P3 & 0 & -- & -- & 88.3 & 156.0 & 161.8 \\
F12\_S07 & central\_dpp & mid & P1 & 30 & 0.8667 & 0.0000 & 2137.4 & 2178.3 & 2188.0 \\
F12\_S07 & central\_dpp & mid & P2 & 30 & 0.8667 & 0.0000 & 861.1 & 927.6 & 931.2 \\
F12\_S07 & central\_dpp & mid & P3 & 30 & 1.0000 & 1.0000 & 87.2 & 159.0 & 164.3 \\
F12\_S07 & central\_dpp & high & P1 & 90 & 0.9778 & 0.0000 & 2061.6 & 2089.0 & 2093.5 \\
F12\_S07 & central\_dpp & high & P2 & 90 & 0.9333 & 0.0000 & 864.5 & 941.0 & 944.8 \\
F12\_S07 & central\_dpp & high & P3 & 90 & 1.0000 & 1.0000 & 81.4 & 144.8 & 150.1 \\
F12\_S07 & shielded\_dpp & low & P1 & 0 & -- & -- & 1894.6 & 1993.4 & 1996.3 \\
F12\_S07 & shielded\_dpp & low & P2 & 0 & -- & -- & 628.4 & 734.7 & 738.1 \\
F12\_S07 & shielded\_dpp & low & P3 & 0 & -- & -- & 85.6 & 158.3 & 163.7 \\
F12\_S07 & shielded\_dpp & mid & P1 & 30 & 0.9333 & 0.0000 & 2068.2 & 2134.6 & 2137.8 \\
F12\_S07 & shielded\_dpp & mid & P2 & 30 & 0.9333 & 0.0000 & 716.9 & 752.4 & 755.3 \\
F12\_S07 & shielded\_dpp & mid & P3 & 30 & 1.0000 & 1.0000 & 79.0 & 149.3 & 154.3 \\
F12\_S07 & shielded\_dpp & high & P1 & 90 & 0.9889 & 0.0000 & 1891.0 & 1956.7 & 1959.9 \\
F12\_S07 & shielded\_dpp & high & P2 & 90 & 0.9556 & 0.0000 & 704.9 & 781.2 & 833.6 \\
F12\_S07 & shielded\_dpp & high & P3 & 90 & 1.0000 & 1.0000 & 79.4 & 144.9 & 151.5 \\
F12\_S08 & fixed\_local & low & P1 & 0 & -- & -- & 1957.4 & 2013.6 & 2016.1 \\
F12\_S08 & fixed\_local & low & P2 & 0 & -- & -- & 685.9 & 732.8 & 735.9 \\
F12\_S08 & fixed\_local & low & P3 & 0 & -- & -- & 87.0 & 158.9 & 165.4 \\
F12\_S08 & fixed\_local & mid & P1 & 0 & -- & -- & 1817.2 & 1937.7 & 1954.1 \\
F12\_S08 & fixed\_local & mid & P2 & 0 & -- & -- & 614.7 & 724.4 & 744.3 \\
F12\_S08 & fixed\_local & mid & P3 & 0 & -- & -- & 88.0 & 158.5 & 164.9 \\
F12\_S08 & fixed\_local & high & P1 & 0 & -- & -- & 1937.6 & 2017.9 & 2035.3 \\
F12\_S08 & fixed\_local & high & P2 & 0 & -- & -- & 683.8 & 776.3 & 784.7 \\
F12\_S08 & fixed\_local & high & P3 & 0 & -- & -- & 89.9 & 164.2 & 170.0 \\
F12\_S08 & fixed\_summary & low & P1 & 0 & -- & -- & 1999.1 & 2079.8 & 2083.7 \\
F12\_S08 & fixed\_summary & low & P2 & 0 & -- & -- & 735.8 & 778.9 & 786.1 \\
F12\_S08 & fixed\_summary & low & P3 & 0 & -- & -- & 88.0 & 158.4 & 164.5 \\
F12\_S08 & fixed\_summary & mid & P1 & 0 & -- & -- & 1992.5 & 2077.9 & 2082.7 \\
F12\_S08 & fixed\_summary & mid & P2 & 0 & -- & -- & 745.7 & 848.4 & 854.3 \\
F12\_S08 & fixed\_summary & mid & P3 & 0 & -- & -- & 89.3 & 158.2 & 164.7 \\
F12\_S08 & fixed\_summary & high & P1 & 0 & -- & -- & 2038.2 & 2116.8 & 2119.5 \\
F12\_S08 & fixed\_summary & high & P2 & 0 & -- & -- & 750.9 & 807.1 & 811.8 \\
F12\_S08 & fixed\_summary & high & P3 & 0 & -- & -- & 87.3 & 161.2 & 167.8 \\
F12\_S08 & fixed\_full & low & P1 & 90 & 0.0000 & 0.0000 & 2158.3 & 2203.4 & 2206.0 \\
F12\_S08 & fixed\_full & low & P2 & 90 & 0.0000 & 0.0000 & 978.1 & 1030.7 & 1034.2 \\
F12\_S08 & fixed\_full & low & P3 & 90 & 0.0000 & 0.0000 & 75.1 & 136.5 & 142.3 \\
F12\_S08 & fixed\_full & mid & P1 & 90 & 0.0000 & 0.0000 & 1956.9 & 2041.8 & 2044.2 \\
F12\_S08 & fixed\_full & mid & P2 & 90 & 0.0000 & 0.0000 & 904.9 & 1014.7 & 1018.5 \\
F12\_S08 & fixed\_full & mid & P3 & 90 & 0.0000 & 0.0000 & 80.4 & 142.8 & 147.6 \\
F12\_S08 & fixed\_full & high & P1 & 90 & 0.9333 & 0.0000 & 2053.0 & 2190.7 & 2193.1 \\
F12\_S08 & fixed\_full & high & P2 & 90 & 0.9556 & 0.0000 & 908.0 & 1033.1 & 1036.6 \\
F12\_S08 & fixed\_full & high & P3 & 90 & 1.0000 & 1.0000 & 80.5 & 142.0 & 146.8 \\
F12\_S08 & independent\_dpp & low & P1 & 22 & 0.0000 & 0.0000 & 1874.7 & 2010.2 & 2012.1 \\
F12\_S08 & independent\_dpp & low & P2 & 22 & 0.0000 & 0.0000 & 773.6 & 889.6 & 892.9 \\
F12\_S08 & independent\_dpp & low & P3 & 22 & 0.0000 & 0.0000 & 82.8 & 150.6 & 158.3 \\
F12\_S08 & independent\_dpp & mid & P1 & 45 & 0.0000 & 0.0000 & 1854.8 & 1911.5 & 1913.6 \\
F12\_S08 & independent\_dpp & mid & P2 & 45 & 0.0000 & 0.0000 & 664.3 & 721.6 & 724.3 \\
F12\_S08 & independent\_dpp & mid & P3 & 45 & 0.0000 & 0.0000 & 78.2 & 140.8 & 145.7 \\
F12\_S08 & independent\_dpp & high & P1 & 90 & 0.9222 & 0.0000 & 2146.9 & 2229.6 & 2231.4 \\
F12\_S08 & independent\_dpp & high & P2 & 90 & 0.9333 & 0.0000 & 995.3 & 1104.8 & 1107.4 \\
F12\_S08 & independent\_dpp & high & P3 & 90 & 1.0000 & 1.0000 & 75.3 & 141.2 & 146.0 \\
F12\_S08 & central\_dpp & low & P1 & 0 & -- & -- & 2023.5 & 2362.2 & 2365.6 \\
F12\_S08 & central\_dpp & low & P2 & 0 & -- & -- & 750.6 & 1083.7 & 1087.1 \\
F12\_S08 & central\_dpp & low & P3 & 0 & -- & -- & 83.0 & 154.3 & 161.6 \\
F12\_S08 & central\_dpp & mid & P1 & 30 & 0.9667 & 0.0000 & 2057.4 & 2105.0 & 2107.2 \\
F12\_S08 & central\_dpp & mid & P2 & 30 & 0.6667 & 0.0000 & 784.2 & 856.1 & 858.2 \\
F12\_S08 & central\_dpp & mid & P3 & 30 & 1.0000 & 1.0000 & 79.3 & 146.0 & 152.4 \\
F12\_S08 & central\_dpp & high & P1 & 90 & 0.9222 & 0.0000 & 2048.7 & 2112.0 & 2115.0 \\
F12\_S08 & central\_dpp & high & P2 & 90 & 0.9556 & 0.0000 & 814.7 & 873.6 & 877.0 \\
F12\_S08 & central\_dpp & high & P3 & 90 & 1.0000 & 1.0000 & 80.4 & 143.1 & 148.8 \\
F12\_S08 & shielded\_dpp & low & P1 & 0 & -- & -- & 1979.4 & 2072.9 & 2075.3 \\
F12\_S08 & shielded\_dpp & low & P2 & 0 & -- & -- & 746.4 & 815.6 & 819.0 \\
F12\_S08 & shielded\_dpp & low & P3 & 0 & -- & -- & 87.2 & 154.6 & 160.7 \\
F12\_S08 & shielded\_dpp & mid & P1 & 30 & 0.7333 & 0.0000 & 1895.4 & 1977.8 & 1980.2 \\
F12\_S08 & shielded\_dpp & mid & P2 & 30 & 0.7000 & 0.0000 & 754.9 & 834.4 & 837.4 \\
F12\_S08 & shielded\_dpp & mid & P3 & 30 & 1.0000 & 1.0000 & 84.4 & 149.1 & 155.6 \\
F12\_S08 & shielded\_dpp & high & P1 & 90 & 0.9556 & 0.0000 & 2070.2 & 2145.7 & 2150.4 \\
F12\_S08 & shielded\_dpp & high & P2 & 90 & 0.9556 & 0.0000 & 946.3 & 1008.9 & 1011.7 \\
F12\_S08 & shielded\_dpp & high & P3 & 90 & 1.0000 & 1.0000 & 72.2 & 135.7 & 140.7 \\
F12\_S09 & fixed\_local & low & P1 & 0 & -- & -- & 2062.2 & 2148.9 & 2153.1 \\
F12\_S09 & fixed\_local & low & P2 & 0 & -- & -- & 835.3 & 908.6 & 912.9 \\
F12\_S09 & fixed\_local & low & P3 & 0 & -- & -- & 91.8 & 166.5 & 174.5 \\
F12\_S09 & fixed\_local & mid & P1 & 0 & -- & -- & 2017.6 & 2068.9 & 2073.1 \\
F12\_S09 & fixed\_local & mid & P2 & 0 & -- & -- & 734.4 & 788.0 & 793.2 \\
F12\_S09 & fixed\_local & mid & P3 & 0 & -- & -- & 84.3 & 154.2 & 160.5 \\
F12\_S09 & fixed\_local & high & P1 & 0 & -- & -- & 2105.4 & 2162.6 & 2166.2 \\
F12\_S09 & fixed\_local & high & P2 & 0 & -- & -- & 862.7 & 933.5 & 937.1 \\
F12\_S09 & fixed\_local & high & P3 & 0 & -- & -- & 90.6 & 162.1 & 167.7 \\
F12\_S09 & fixed\_summary & low & P1 & 0 & -- & -- & 2035.6 & 2133.6 & 2137.1 \\
F12\_S09 & fixed\_summary & low & P2 & 0 & -- & -- & 740.9 & 828.5 & 832.8 \\
F12\_S09 & fixed\_summary & low & P3 & 0 & -- & -- & 91.4 & 166.4 & 172.2 \\
F12\_S09 & fixed\_summary & mid & P1 & 0 & -- & -- & 2067.9 & 2207.0 & 2209.7 \\
F12\_S09 & fixed\_summary & mid & P2 & 0 & -- & -- & 731.5 & 865.4 & 869.0 \\
F12\_S09 & fixed\_summary & mid & P3 & 0 & -- & -- & 89.9 & 163.7 & 169.8 \\
F12\_S09 & fixed\_summary & high & P1 & 0 & -- & -- & 1931.7 & 2076.8 & 2079.4 \\
F12\_S09 & fixed\_summary & high & P2 & 0 & -- & -- & 656.3 & 775.4 & 778.8 \\
F12\_S09 & fixed\_summary & high & P3 & 0 & -- & -- & 84.7 & 152.8 & 158.7 \\
F12\_S09 & fixed\_full & low & P1 & 90 & 0.0000 & 0.0000 & 2058.8 & 2117.3 & 2124.4 \\
F12\_S09 & fixed\_full & low & P2 & 90 & 0.0000 & 0.0000 & 756.7 & 835.5 & 843.0 \\
F12\_S09 & fixed\_full & low & P3 & 90 & 0.0000 & 0.0000 & 77.2 & 141.2 & 148.3 \\
F12\_S09 & fixed\_full & mid & P1 & 90 & 0.0000 & 0.0000 & 2005.3 & 2062.5 & 2089.6 \\
F12\_S09 & fixed\_full & mid & P2 & 90 & 0.0000 & 0.0000 & 877.8 & 931.8 & 935.3 \\
F12\_S09 & fixed\_full & mid & P3 & 90 & 0.0000 & 0.0000 & 77.9 & 136.7 & 142.7 \\
F12\_S09 & fixed\_full & high & P1 & 90 & 0.9111 & 0.0000 & 2041.2 & 2101.3 & 2106.6 \\
F12\_S09 & fixed\_full & high & P2 & 90 & 0.9556 & 0.0000 & 795.5 & 870.5 & 874.5 \\
F12\_S09 & fixed\_full & high & P3 & 90 & 1.0000 & 1.0000 & 77.1 & 141.8 & 147.0 \\
F12\_S09 & independent\_dpp & low & P1 & 22 & 0.0000 & 0.0000 & 2044.8 & 2140.4 & 2142.8 \\
F12\_S09 & independent\_dpp & low & P2 & 22 & 0.0000 & 0.0000 & 886.4 & 937.7 & 959.9 \\
F12\_S09 & independent\_dpp & low & P3 & 22 & 0.0000 & 0.0000 & 85.2 & 148.8 & 157.8 \\
F12\_S09 & independent\_dpp & mid & P1 & 45 & 0.0000 & 0.0000 & 2163.0 & 2195.9 & 2197.6 \\
F12\_S09 & independent\_dpp & mid & P2 & 45 & 0.0000 & 0.0000 & 1015.9 & 1045.1 & 1065.2 \\
F12\_S09 & independent\_dpp & mid & P3 & 45 & 0.0000 & 0.0000 & 78.9 & 140.5 & 145.3 \\
F12\_S09 & independent\_dpp & high & P1 & 90 & 0.9333 & 0.0000 & 2125.5 & 2172.7 & 2174.9 \\
F12\_S09 & independent\_dpp & high & P2 & 90 & 0.9333 & 0.0000 & 851.2 & 955.2 & 959.5 \\
F12\_S09 & independent\_dpp & high & P3 & 90 & 1.0000 & 1.0000 & 74.2 & 135.0 & 139.5 \\
F12\_S09 & central\_dpp & low & P1 & 0 & -- & -- & 1974.6 & 2062.7 & 2073.2 \\
F12\_S09 & central\_dpp & low & P2 & 0 & -- & -- & 700.8 & 778.4 & 783.7 \\
F12\_S09 & central\_dpp & low & P3 & 0 & -- & -- & 88.6 & 160.8 & 166.0 \\
F12\_S09 & central\_dpp & mid & P1 & 30 & 0.9000 & 0.0000 & 2105.6 & 2140.7 & 2143.5 \\
F12\_S09 & central\_dpp & mid & P2 & 30 & 0.7667 & 0.0000 & 957.1 & 996.2 & 1018.1 \\
F12\_S09 & central\_dpp & mid & P3 & 30 & 1.0000 & 1.0000 & 85.9 & 153.4 & 158.5 \\
F12\_S09 & central\_dpp & high & P1 & 90 & 0.9556 & 0.0000 & 2328.7 & 2368.8 & 2373.6 \\
F12\_S09 & central\_dpp & high & P2 & 90 & 0.9778 & 0.0000 & 997.1 & 1036.6 & 1039.8 \\
F12\_S09 & central\_dpp & high & P3 & 90 & 1.0000 & 1.0000 & 68.1 & 135.1 & 142.1 \\
F12\_S09 & shielded\_dpp & low & P1 & 0 & -- & -- & 1968.5 & 2040.6 & 2046.2 \\
F12\_S09 & shielded\_dpp & low & P2 & 0 & -- & -- & 790.4 & 868.6 & 870.5 \\
F12\_S09 & shielded\_dpp & low & P3 & 0 & -- & -- & 92.7 & 162.5 & 168.8 \\
F12\_S09 & shielded\_dpp & mid & P1 & 30 & 0.8667 & 0.0000 & 2063.5 & 2130.1 & 2140.9 \\
F12\_S09 & shielded\_dpp & mid & P2 & 30 & 0.9000 & 0.0000 & 834.2 & 885.0 & 888.3 \\
F12\_S09 & shielded\_dpp & mid & P3 & 30 & 1.0000 & 1.0000 & 80.4 & 145.5 & 150.6 \\
F12\_S09 & shielded\_dpp & high & P1 & 90 & 0.9444 & 0.0000 & 2087.2 & 2138.8 & 2140.9 \\
F12\_S09 & shielded\_dpp & high & P2 & 90 & 0.9444 & 0.0000 & 824.1 & 897.0 & 899.7 \\
F12\_S09 & shielded\_dpp & high & P3 & 90 & 1.0000 & 1.0000 & 73.9 & 137.8 & 143.0 \\
F12\_S10 & fixed\_local & low & P1 & 0 & -- & -- & 1963.5 & 2064.8 & 2068.2 \\
F12\_S10 & fixed\_local & low & P2 & 0 & -- & -- & 761.3 & 837.0 & 839.8 \\
F12\_S10 & fixed\_local & low & P3 & 0 & -- & -- & 91.5 & 166.4 & 173.2 \\
F12\_S10 & fixed\_local & mid & P1 & 0 & -- & -- & 2011.0 & 2050.3 & 2053.7 \\
F12\_S10 & fixed\_local & mid & P2 & 0 & -- & -- & 715.1 & 755.4 & 758.7 \\
F12\_S10 & fixed\_local & mid & P3 & 0 & -- & -- & 90.9 & 162.7 & 169.3 \\
F12\_S10 & fixed\_local & high & P1 & 0 & -- & -- & 1937.4 & 2000.6 & 2003.6 \\
F12\_S10 & fixed\_local & high & P2 & 0 & -- & -- & 668.6 & 723.4 & 726.7 \\
F12\_S10 & fixed\_local & high & P3 & 0 & -- & -- & 93.9 & 163.6 & 170.2 \\
F12\_S10 & fixed\_summary & low & P1 & 0 & -- & -- & 2013.5 & 2103.8 & 2108.3 \\
F12\_S10 & fixed\_summary & low & P2 & 0 & -- & -- & 823.8 & 913.5 & 917.0 \\
F12\_S10 & fixed\_summary & low & P3 & 0 & -- & -- & 84.6 & 153.6 & 160.5 \\
F12\_S10 & fixed\_summary & mid & P1 & 0 & -- & -- & 2057.6 & 2136.9 & 2140.5 \\
F12\_S10 & fixed\_summary & mid & P2 & 0 & -- & -- & 946.8 & 1014.3 & 1018.1 \\
F12\_S10 & fixed\_summary & mid & P3 & 0 & -- & -- & 84.4 & 155.9 & 163.6 \\
F12\_S10 & fixed\_summary & high & P1 & 0 & -- & -- & 1943.4 & 2129.7 & 2134.3 \\
F12\_S10 & fixed\_summary & high & P2 & 0 & -- & -- & 764.8 & 939.7 & 943.1 \\
F12\_S10 & fixed\_summary & high & P3 & 0 & -- & -- & 87.3 & 163.1 & 170.0 \\
F12\_S10 & fixed\_full & low & P1 & 90 & 0.0000 & 0.0000 & 1967.7 & 2037.9 & 2041.3 \\
F12\_S10 & fixed\_full & low & P2 & 90 & 0.0000 & 0.0000 & 858.5 & 951.5 & 953.8 \\
F12\_S10 & fixed\_full & low & P3 & 90 & 0.0000 & 0.0000 & 76.8 & 138.7 & 143.3 \\
F12\_S10 & fixed\_full & mid & P1 & 90 & 0.0000 & 0.0000 & 2146.3 & 2221.4 & 2223.5 \\
F12\_S10 & fixed\_full & mid & P2 & 90 & 0.0000 & 0.0000 & 847.3 & 890.1 & 894.4 \\
F12\_S10 & fixed\_full & mid & P3 & 90 & 0.0000 & 0.0000 & 74.8 & 132.3 & 136.9 \\
F12\_S10 & fixed\_full & high & P1 & 90 & 0.9778 & 0.0000 & 2165.6 & 2199.9 & 2204.4 \\
F12\_S10 & fixed\_full & high & P2 & 90 & 0.9667 & 0.0000 & 871.9 & 912.1 & 916.5 \\
F12\_S10 & fixed\_full & high & P3 & 90 & 1.0000 & 1.0000 & 81.3 & 142.9 & 147.6 \\
F12\_S10 & independent\_dpp & low & P1 & 22 & 0.0000 & 0.0000 & 1921.7 & 2011.9 & 2016.7 \\
F12\_S10 & independent\_dpp & low & P2 & 22 & 0.0000 & 0.0000 & 653.8 & 788.2 & 791.3 \\
F12\_S10 & independent\_dpp & low & P3 & 22 & 0.0000 & 0.0000 & 82.9 & 150.7 & 156.6 \\
F12\_S10 & independent\_dpp & mid & P1 & 45 & 0.0000 & 0.0000 & 1887.2 & 1966.4 & 1968.3 \\
F12\_S10 & independent\_dpp & mid & P2 & 45 & 0.0000 & 0.0000 & 773.0 & 838.0 & 840.0 \\
F12\_S10 & independent\_dpp & mid & P3 & 45 & 0.0000 & 0.0000 & 75.8 & 141.7 & 147.3 \\
F12\_S10 & independent\_dpp & high & P1 & 90 & 0.8556 & 0.0000 & 1965.1 & 2090.0 & 2093.2 \\
F12\_S10 & independent\_dpp & high & P2 & 90 & 0.9667 & 0.0000 & 882.8 & 981.7 & 984.1 \\
F12\_S10 & independent\_dpp & high & P3 & 90 & 1.0000 & 1.0000 & 76.0 & 142.0 & 147.0 \\
F12\_S10 & central\_dpp & low & P1 & 0 & -- & -- & 2041.7 & 2115.9 & 2119.5 \\
F12\_S10 & central\_dpp & low & P2 & 0 & -- & -- & 834.2 & 905.3 & 909.7 \\
F12\_S10 & central\_dpp & low & P3 & 0 & -- & -- & 89.3 & 164.3 & 169.0 \\
F12\_S10 & central\_dpp & mid & P1 & 30 & 0.9667 & 0.0000 & 2076.9 & 2105.3 & 2151.2 \\
F12\_S10 & central\_dpp & mid & P2 & 30 & 0.9000 & 0.0000 & 850.4 & 929.6 & 931.8 \\
F12\_S10 & central\_dpp & mid & P3 & 30 & 1.0000 & 1.0000 & 80.3 & 148.5 & 154.8 \\
F12\_S10 & central\_dpp & high & P1 & 90 & 0.9556 & 0.0000 & 2040.7 & 2102.1 & 2106.1 \\
F12\_S10 & central\_dpp & high & P2 & 90 & 0.9667 & 0.0000 & 905.9 & 940.3 & 944.7 \\
F12\_S10 & central\_dpp & high & P3 & 90 & 1.0000 & 1.0000 & 76.4 & 140.4 & 146.7 \\
F12\_S10 & shielded\_dpp & low & P1 & 0 & -- & -- & 2085.7 & 2128.5 & 2130.9 \\
F12\_S10 & shielded\_dpp & low & P2 & 0 & -- & -- & 938.2 & 977.5 & 981.4 \\
F12\_S10 & shielded\_dpp & low & P3 & 0 & -- & -- & 89.8 & 160.9 & 166.9 \\
F12\_S10 & shielded\_dpp & mid & P1 & 30 & 0.8000 & 0.0000 & 1980.5 & 2053.7 & 2058.1 \\
F12\_S10 & shielded\_dpp & mid & P2 & 30 & 0.7667 & 0.0000 & 729.0 & 827.8 & 830.2 \\
F12\_S10 & shielded\_dpp & mid & P3 & 30 & 1.0000 & 1.0000 & 86.9 & 154.8 & 160.6 \\
F12\_S10 & shielded\_dpp & high & P1 & 90 & 0.9222 & 0.0000 & 2020.4 & 2096.1 & 2123.6 \\
F12\_S10 & shielded\_dpp & high & P2 & 90 & 0.9778 & 0.0000 & 879.9 & 928.6 & 932.5 \\
F12\_S10 & shielded\_dpp & high & P3 & 90 & 1.0000 & 1.0000 & 80.1 & 143.5 & 149.0 \\
F12\_S11 & fixed\_local & low & P1 & 0 & -- & -- & 1923.5 & 1986.9 & 2001.2 \\
F12\_S11 & fixed\_local & low & P2 & 0 & -- & -- & 743.9 & 820.1 & 835.6 \\
F12\_S11 & fixed\_local & low & P3 & 0 & -- & -- & 89.1 & 161.4 & 168.1 \\
F12\_S11 & fixed\_local & mid & P1 & 0 & -- & -- & 2032.4 & 2109.7 & 2112.1 \\
F12\_S11 & fixed\_local & mid & P2 & 0 & -- & -- & 748.3 & 824.6 & 828.8 \\
F12\_S11 & fixed\_local & mid & P3 & 0 & -- & -- & 93.9 & 169.5 & 176.2 \\
F12\_S11 & fixed\_local & high & P1 & 0 & -- & -- & 2031.0 & 2105.3 & 2109.9 \\
F12\_S11 & fixed\_local & high & P2 & 0 & -- & -- & 736.1 & 814.3 & 818.7 \\
F12\_S11 & fixed\_local & high & P3 & 0 & -- & -- & 90.1 & 162.4 & 168.3 \\
F12\_S11 & fixed\_summary & low & P1 & 0 & -- & -- & 1928.4 & 1983.9 & 1990.1 \\
F12\_S11 & fixed\_summary & low & P2 & 0 & -- & -- & 639.4 & 714.9 & 725.7 \\
F12\_S11 & fixed\_summary & low & P3 & 0 & -- & -- & 82.4 & 153.2 & 159.2 \\
F12\_S11 & fixed\_summary & mid & P1 & 0 & -- & -- & 2087.9 & 2136.3 & 2140.6 \\
F12\_S11 & fixed\_summary & mid & P2 & 0 & -- & -- & 824.7 & 887.7 & 894.9 \\
F12\_S11 & fixed\_summary & mid & P3 & 0 & -- & -- & 86.8 & 157.3 & 162.5 \\
F12\_S11 & fixed\_summary & high & P1 & 0 & -- & -- & 1941.6 & 1992.7 & 1995.2 \\
F12\_S11 & fixed\_summary & high & P2 & 0 & -- & -- & 648.9 & 703.8 & 707.1 \\
F12\_S11 & fixed\_summary & high & P3 & 0 & -- & -- & 86.9 & 155.8 & 163.5 \\
F12\_S11 & fixed\_full & low & P1 & 90 & 0.0000 & 0.0000 & 1942.0 & 2002.2 & 2006.4 \\
F12\_S11 & fixed\_full & low & P2 & 90 & 0.0000 & 0.0000 & 724.1 & 806.1 & 808.5 \\
F12\_S11 & fixed\_full & low & P3 & 90 & 0.0000 & 0.0000 & 85.2 & 150.6 & 155.4 \\
F12\_S11 & fixed\_full & mid & P1 & 90 & 0.0000 & 0.0000 & 2070.3 & 2149.2 & 2152.6 \\
F12\_S11 & fixed\_full & mid & P2 & 90 & 0.0000 & 0.0000 & 938.8 & 980.4 & 983.0 \\
F12\_S11 & fixed\_full & mid & P3 & 90 & 0.0000 & 0.0000 & 79.7 & 142.6 & 148.2 \\
F12\_S11 & fixed\_full & high & P1 & 90 & 0.9111 & 0.0000 & 2113.8 & 2201.4 & 2203.3 \\
F12\_S11 & fixed\_full & high & P2 & 90 & 0.9333 & 0.0000 & 912.5 & 1038.8 & 1042.4 \\
F12\_S11 & fixed\_full & high & P3 & 90 & 1.0000 & 1.0000 & 76.8 & 135.8 & 141.6 \\
F12\_S11 & independent\_dpp & low & P1 & 22 & 0.0000 & 0.0000 & 1986.8 & 2019.7 & 2022.3 \\
F12\_S11 & independent\_dpp & low & P2 & 22 & 0.0000 & 0.0000 & 764.0 & 802.3 & 804.8 \\
F12\_S11 & independent\_dpp & low & P3 & 22 & 0.0000 & 0.0000 & 88.9 & 156.0 & 162.7 \\
F12\_S11 & independent\_dpp & mid & P1 & 45 & 0.0000 & 0.0000 & 2001.6 & 2087.2 & 2091.6 \\
F12\_S11 & independent\_dpp & mid & P2 & 45 & 0.0000 & 0.0000 & 824.6 & 904.2 & 906.9 \\
F12\_S11 & independent\_dpp & mid & P3 & 45 & 0.0000 & 0.0000 & 78.9 & 141.5 & 146.5 \\
F12\_S11 & independent\_dpp & high & P1 & 90 & 0.9000 & 0.0000 & 2133.6 & 2182.2 & 2184.9 \\
F12\_S11 & independent\_dpp & high & P2 & 90 & 0.9444 & 0.0000 & 981.2 & 1046.9 & 1051.6 \\
F12\_S11 & independent\_dpp & high & P3 & 90 & 1.0000 & 1.0000 & 71.8 & 129.6 & 136.5 \\
F12\_S11 & central\_dpp & low & P1 & 0 & -- & -- & 2043.7 & 2116.0 & 2120.4 \\
F12\_S11 & central\_dpp & low & P2 & 0 & -- & -- & 783.2 & 865.3 & 868.8 \\
F12\_S11 & central\_dpp & low & P3 & 0 & -- & -- & 93.5 & 162.4 & 168.8 \\
F12\_S11 & central\_dpp & mid & P1 & 30 & 0.8000 & 0.0000 & 2106.5 & 2167.1 & 2169.4 \\
F12\_S11 & central\_dpp & mid & P2 & 30 & 0.8333 & 0.0000 & 934.4 & 991.7 & 995.4 \\
F12\_S11 & central\_dpp & mid & P3 & 30 & 1.0000 & 1.0000 & 80.0 & 146.7 & 151.1 \\
F12\_S11 & central\_dpp & high & P1 & 90 & 0.8889 & 0.0000 & 1923.1 & 1979.1 & 1982.6 \\
F12\_S11 & central\_dpp & high & P2 & 90 & 0.9444 & 0.0000 & 689.4 & 782.6 & 785.7 \\
F12\_S11 & central\_dpp & high & P3 & 90 & 1.0000 & 1.0000 & 75.7 & 135.1 & 140.8 \\
F12\_S11 & shielded\_dpp & low & P1 & 0 & -- & -- & 2014.5 & 2077.7 & 2080.4 \\
F12\_S11 & shielded\_dpp & low & P2 & 0 & -- & -- & 735.4 & 818.5 & 822.9 \\
F12\_S11 & shielded\_dpp & low & P3 & 0 & -- & -- & 87.8 & 158.6 & 165.2 \\
F12\_S11 & shielded\_dpp & mid & P1 & 30 & 0.8333 & 0.0000 & 1884.2 & 1947.6 & 1950.1 \\
F12\_S11 & shielded\_dpp & mid & P2 & 30 & 0.8667 & 0.0000 & 802.2 & 834.0 & 882.2 \\
F12\_S11 & shielded\_dpp & mid & P3 & 30 & 1.0000 & 1.0000 & 79.7 & 142.8 & 149.4 \\
F12\_S11 & shielded\_dpp & high & P1 & 90 & 0.9000 & 0.0000 & 1866.9 & 1926.8 & 1928.9 \\
F12\_S11 & shielded\_dpp & high & P2 & 90 & 0.9333 & 0.0000 & 651.9 & 721.0 & 724.2 \\
F12\_S11 & shielded\_dpp & high & P3 & 90 & 1.0000 & 1.0000 & 78.9 & 143.3 & 149.7 \\
F12\_S12 & fixed\_local & low & P1 & 0 & -- & -- & 2021.6 & 2060.8 & 2064.9 \\
F12\_S12 & fixed\_local & low & P2 & 0 & -- & -- & 851.3 & 891.4 & 896.0 \\
F12\_S12 & fixed\_local & low & P3 & 0 & -- & -- & 89.4 & 162.6 & 169.8 \\
F12\_S12 & fixed\_local & mid & P1 & 0 & -- & -- & 2065.4 & 2110.6 & 2113.4 \\
F12\_S12 & fixed\_local & mid & P2 & 0 & -- & -- & 880.8 & 918.2 & 921.2 \\
F12\_S12 & fixed\_local & mid & P3 & 0 & -- & -- & 95.8 & 167.4 & 173.3 \\
F12\_S12 & fixed\_local & high & P1 & 0 & -- & -- & 2028.6 & 2158.6 & 2162.3 \\
F12\_S12 & fixed\_local & high & P2 & 0 & -- & -- & 775.1 & 837.5 & 843.8 \\
F12\_S12 & fixed\_local & high & P3 & 0 & -- & -- & 92.3 & 164.3 & 171.5 \\
F12\_S12 & fixed\_summary & low & P1 & 0 & -- & -- & 1910.3 & 1985.2 & 1989.2 \\
F12\_S12 & fixed\_summary & low & P2 & 0 & -- & -- & 718.9 & 806.6 & 808.8 \\
F12\_S12 & fixed\_summary & low & P3 & 0 & -- & -- & 88.1 & 159.9 & 165.1 \\
F12\_S12 & fixed\_summary & mid & P1 & 0 & -- & -- & 1893.9 & 1959.5 & 1962.4 \\
F12\_S12 & fixed\_summary & mid & P2 & 0 & -- & -- & 613.9 & 675.0 & 678.6 \\
F12\_S12 & fixed\_summary & mid & P3 & 0 & -- & -- & 88.1 & 155.1 & 161.6 \\
F12\_S12 & fixed\_summary & high & P1 & 0 & -- & -- & 2217.8 & 2256.9 & 2260.5 \\
F12\_S12 & fixed\_summary & high & P2 & 0 & -- & -- & 1057.1 & 1098.5 & 1102.7 \\
F12\_S12 & fixed\_summary & high & P3 & 0 & -- & -- & 82.4 & 150.6 & 156.6 \\
F12\_S12 & fixed\_full & low & P1 & 90 & 0.0000 & 0.0000 & 2056.0 & 2127.3 & 2128.7 \\
F12\_S12 & fixed\_full & low & P2 & 90 & 0.0000 & 0.0000 & 831.5 & 884.5 & 889.1 \\
F12\_S12 & fixed\_full & low & P3 & 90 & 0.0000 & 0.0000 & 85.1 & 157.7 & 163.5 \\
F12\_S12 & fixed\_full & mid & P1 & 90 & 0.0000 & 0.0000 & 1859.9 & 1997.8 & 2002.3 \\
F12\_S12 & fixed\_full & mid & P2 & 90 & 0.0000 & 0.0000 & 697.3 & 841.4 & 844.7 \\
F12\_S12 & fixed\_full & mid & P3 & 90 & 0.0000 & 0.0000 & 83.0 & 145.3 & 151.3 \\
F12\_S12 & fixed\_full & high & P1 & 90 & 0.9222 & 0.0000 & 2070.3 & 2135.0 & 2139.0 \\
F12\_S12 & fixed\_full & high & P2 & 90 & 0.9444 & 0.0000 & 956.1 & 1001.6 & 1003.5 \\
F12\_S12 & fixed\_full & high & P3 & 90 & 1.0000 & 1.0000 & 80.0 & 138.6 & 145.2 \\
F12\_S12 & independent\_dpp & low & P1 & 22 & 0.0000 & 0.0000 & 1001.7 & 1038.7 & 1041.5 \\
F12\_S12 & independent\_dpp & low & P2 & 22 & 0.0000 & 0.0000 & 864.1 & 915.7 & 918.1 \\
F12\_S12 & independent\_dpp & low & P3 & 22 & 0.0000 & 0.0000 & 85.0 & 157.5 & 162.6 \\
F12\_S12 & independent\_dpp & mid & P1 & 45 & 0.0000 & 0.0000 & 2161.5 & 2193.0 & 2197.7 \\
F12\_S12 & independent\_dpp & mid & P2 & 45 & 0.0000 & 0.0000 & 860.3 & 896.0 & 897.8 \\
F12\_S12 & independent\_dpp & mid & P3 & 45 & 0.0000 & 0.0000 & 77.0 & 138.3 & 143.6 \\
F12\_S12 & independent\_dpp & high & P1 & 90 & 0.9444 & 0.0000 & 2004.9 & 2052.8 & 2056.1 \\
F12\_S12 & independent\_dpp & high & P2 & 90 & 0.9444 & 0.0000 & 778.7 & 819.1 & 822.4 \\
F12\_S12 & independent\_dpp & high & P3 & 90 & 1.0000 & 1.0000 & 74.8 & 135.8 & 141.5 \\
F12\_S12 & central\_dpp & low & P1 & 0 & -- & -- & 2072.6 & 2369.1 & 2372.7 \\
F12\_S12 & central\_dpp & low & P2 & 0 & -- & -- & 865.3 & 1186.7 & 1191.2 \\
F12\_S12 & central\_dpp & low & P3 & 0 & -- & -- & 86.1 & 161.1 & 167.6 \\
F12\_S12 & central\_dpp & mid & P1 & 30 & 0.5333 & 0.0000 & 1883.9 & 1935.5 & 1939.8 \\
F12\_S12 & central\_dpp & mid & P2 & 30 & 0.7667 & 0.0000 & 659.2 & 732.8 & 735.7 \\
F12\_S12 & central\_dpp & mid & P3 & 30 & 1.0000 & 1.0000 & 83.3 & 144.6 & 150.7 \\
F12\_S12 & central\_dpp & high & P1 & 90 & 0.9222 & 0.0000 & 2095.8 & 2187.1 & 2190.5 \\
F12\_S12 & central\_dpp & high & P2 & 90 & 0.9444 & 0.0000 & 846.9 & 919.8 & 922.5 \\
F12\_S12 & central\_dpp & high & P3 & 90 & 1.0000 & 1.0000 & 81.8 & 147.6 & 152.4 \\
F12\_S12 & shielded\_dpp & low & P1 & 0 & -- & -- & 1879.5 & 1960.7 & 1964.6 \\
F12\_S12 & shielded\_dpp & low & P2 & 0 & -- & -- & 695.3 & 756.2 & 759.7 \\
F12\_S12 & shielded\_dpp & low & P3 & 0 & -- & -- & 87.3 & 156.0 & 161.3 \\
F12\_S12 & shielded\_dpp & mid & P1 & 30 & 0.9000 & 0.0000 & 2041.9 & 2073.2 & 2075.5 \\
F12\_S12 & shielded\_dpp & mid & P2 & 30 & 0.9333 & 0.0000 & 803.7 & 835.2 & 837.5 \\
F12\_S12 & shielded\_dpp & mid & P3 & 30 & 1.0000 & 1.0000 & 83.1 & 147.1 & 152.7 \\
F12\_S12 & shielded\_dpp & high & P1 & 90 & 0.9444 & 0.0000 & 1962.8 & 2023.5 & 2026.7 \\
F12\_S12 & shielded\_dpp & high & P2 & 90 & 0.9444 & 0.0000 & 736.5 & 774.2 & 776.1 \\
F12\_S12 & shielded\_dpp & high & P3 & 90 & 1.0000 & 1.0000 & 74.7 & 136.9 & 142.8 \\
\bottomrule
\end{longtable}
\end{landscape}

\begin{landscape}
\section{Release-Lateness Diagnostic}
Independent nominal clocks, the three-probe barrier, unequal first-iteration waiting, relative sleep, and no absolute catch-up are source/log supported. The specific OS/network reason for the recurring proposal-arrival order is UNKNOWN.
\small\begin{longtable}{lrrrrp{65mm}}
\caption{Release-lateness diagnostic from accepted formal timestamps (ms).}\label{tab:release}\\
\toprule
Device & Events & p50 & p95 & p99 & Unproven residual \\
\midrule
\endfirsthead
\toprule
Device & Events & p50 & p95 & p99 & Unproven residual \\
\midrule
\endhead
P1 & 19440 & 1972.34 & 2173.36 & 2326.13 & The exact OS/network reason that proposals repeatedly reached the barrier in P3/P2/P1 order, and residual per-session variation, are UNKNOWN. \\
P2 & 19440 & 763.90 & 976.78 & 1080.30 & The exact OS/network reason that proposals repeatedly reached the barrier in P3/P2/P1 order, and residual per-session variation, are UNKNOWN. \\
P3 & 19440 & 83.56 & 154.08 & 167.94 & The exact OS/network reason that proposals repeatedly reached the barrier in P3/P2/P1 order, and residual per-session variation, are UNKNOWN. \\
\bottomrule
\end{longtable}
\end{landscape}

\begin{landscape}
\section{Finite-n Scaling Conformance}
This is discrete-event conformance, not a 32-probe hardware scalability experiment.
\small\begin{longtable}{rrrrrrrr}
\caption{All finite-$n$ discrete-event conformance cells.}\label{tab:scalefull}\\
\toprule
$n$ & $\rho$ & $C$ (B) & $k^*$ & Sim. $k$ & Theory frac. & Sim. frac. & Match \\
\midrule
\endfirsthead
\toprule
$n$ & $\rho$ & $C$ (B) & $k^*$ & Sim. $k$ & Theory frac. & Sim. frac. & Match \\
\midrule
\endhead
3 & 1/3 & 131072 & 0 & 0 & 0.000000 & 0.000000 & PASS \\
3 & 1/2 & 196608 & 1 & 1 & 0.333333 & 0.333333 & PASS \\
3 & 2/3 & 262144 & 1 & 1 & 0.333333 & 0.333333 & PASS \\
3 & 3/4 & 294912 & 2 & 2 & 0.666667 & 0.666667 & PASS \\
3 & 1/1 & 393216 & 3 & 3 & 1.000000 & 1.000000 & PASS \\
8 & 1/3 & 349525 & 1 & 1 & 0.125000 & 0.125000 & PASS \\
8 & 1/2 & 524288 & 3 & 3 & 0.375000 & 0.375000 & PASS \\
8 & 2/3 & 699050 & 4 & 4 & 0.500000 & 0.500000 & PASS \\
8 & 3/4 & 786432 & 5 & 5 & 0.625000 & 0.625000 & PASS \\
8 & 1/1 & 1048576 & 8 & 8 & 1.000000 & 1.000000 & PASS \\
16 & 1/3 & 699050 & 3 & 3 & 0.187500 & 0.187500 & PASS \\
16 & 1/2 & 1048576 & 6 & 6 & 0.375000 & 0.375000 & PASS \\
16 & 2/3 & 1398101 & 9 & 9 & 0.562500 & 0.562500 & PASS \\
16 & 3/4 & 1572864 & 11 & 11 & 0.687500 & 0.687500 & PASS \\
16 & 1/1 & 2097152 & 16 & 16 & 1.000000 & 1.000000 & PASS \\
32 & 1/3 & 1398101 & 7 & 7 & 0.218750 & 0.218750 & PASS \\
32 & 1/2 & 2097152 & 13 & 13 & 0.406250 & 0.406250 & PASS \\
32 & 2/3 & 2796202 & 19 & 19 & 0.593750 & 0.593750 & PASS \\
32 & 3/4 & 3145728 & 22 & 22 & 0.687500 & 0.687500 & PASS \\
32 & 1/1 & 4194304 & 32 & 32 & 1.000000 & 1.000000 & PASS \\
\bottomrule
\end{longtable}
\end{landscape}

\section{Accepted and Excluded Attempts}
Accepted attempts are S01-a01, S02-a01, S03-a02, S04-a02, S05-a01, S06-a01, S07-a02, and S08--S12-a01. S03-a01, S04-a01, and S07-a01 are excluded infrastructure attempts and are not stitched into accepted sessions.

\section{Detector and Replay Limitations}
The workload is frozen CICIoT2023-derived feature-vector replay, not packet replay or an uncontrolled deployment. The formal occurrence set excludes Recon-PortScan, and exact feature-vector overlap exists across frozen partitions. Perfect correctness on the formal occurrences is therefore a controlled diagnostic rather than an IDS-superiority result.

\section{Raw Audit and Methodology-Repair Verdicts}
The full-raw audit rebuilt 216/216 accepted arms and 58,320 measured events with zero raw-to-metric mismatch and zero shared-cap violation across 63,720 service rounds. Its verdict is PASS\_WITH\_ISSUES because the raw bundle lacks direct artifacts for the S03-a01 SSD incident and the S07 pre-hardware launcher abort. The methodology repair is PASS with retained rows A/B/C/D = 432/240/648/20 and zero failed runs.

\section{Artifact Access}
Machine-readable CSV/JSON sources, number provenance, claim audit, build report, validation report, and SHA-256 manifest accompany this supplementary PDF. Third-party dataset redistribution rights remain outside the scope of this package.
\end{document}
